\documentclass[a4paper, 11pt]{article}

% -- Encoding & language ----------------------------------------------
\usepackage[utf8]{inputenc}
\usepackage[T2A]{fontenc}
\usepackage[russian]{babel}

% -- Maths ------------------------------------------------------------
\usepackage{amsmath, amssymb, amsthm}

% -- Layout -----------------------------------------------------------
\usepackage[margin=2.5cm]{geometry}
\usepackage{parskip}
\usepackage{enumitem}
\usepackage{booktabs}
\usepackage{longtable}
\usepackage{array}
\usepackage{multicol}

% -- Graphics & links -------------------------------------------------
\usepackage{xcolor}
\usepackage[colorlinks=true, linkcolor=blue!70!black, citecolor=green!50!black, urlcolor=blue!60!black]{hyperref}

% -- Bibliography -----------------------------------------------------
\usepackage[numbers]{natbib}

% -- Appendix ---------------------------------------------------------
\usepackage[toc,page]{appendix}

% -- Verbatim for prompts ---------------------------------------------
\usepackage{fancyvrb}
\usepackage{listings}
\lstset{
  basicstyle=\ttfamily\footnotesize,
  breaklines=true,
  frame=single,
  framesep=4pt,
  xleftmargin=8pt,
  xrightmargin=8pt,
  inputencoding=utf8,
  extendedchars=true,
  literate={а}{{\cyra}}1 {б}{{\cyrb}}1 {в}{{\cyrv}}1 {г}{{\cyrg}}1
           {д}{{\cyrd}}1 {е}{{\cyre}}1 {ж}{{\cyrzh}}1 {з}{{\cyrz}}1
           {и}{{\cyri}}1 {й}{{\cyrishrt}}1 {к}{{\cyrk}}1 {л}{{\cyrl}}1
           {м}{{\cyrm}}1 {н}{{\cyrn}}1 {о}{{\cyro}}1 {п}{{\cyrp}}1
           {р}{{\cyrr}}1 {с}{{\cyrs}}1 {т}{{\cyrt}}1 {у}{{\cyru}}1
           {ф}{{\cyrf}}1 {х}{{\cyrh}}1 {ц}{{\cyrc}}1 {ч}{{\cyrch}}1
           {ш}{{\cyrsh}}1 {щ}{{\cyrshch}}1 {ъ}{{\cyrhrdsn}}1 {ы}{{\cyry}}1
           {ь}{{\cyrsftsn}}1 {э}{{\cyrerev}}1 {ю}{{\cyryu}}1 {я}{{\cyrya}}1
           {А}{{\CYRA}}1 {Б}{{\CYRB}}1 {В}{{\CYRV}}1 {Г}{{\CYRG}}1
           {Д}{{\CYRD}}1 {Е}{{\CYRE}}1 {Ж}{{\CYRZH}}1 {З}{{\CYRZ}}1
           {И}{{\CYRI}}1 {Й}{{\CYRISHRT}}1 {К}{{\CYRK}}1 {Л}{{\CYRL}}1
           {М}{{\CYRM}}1 {Н}{{\CYRN}}1 {О}{{\CYRO}}1 {П}{{\CYRP}}1
           {Р}{{\CYRR}}1 {С}{{\CYRS}}1 {Т}{{\CYRT}}1 {У}{{\CYRU}}1
           {Ф}{{\CYRF}}1 {Х}{{\CYRH}}1 {Ц}{{\CYRC}}1 {Ч}{{\CYRCH}}1
           {Ш}{{\CYRSH}}1 {Щ}{{\CYRSHCH}}1 {Ъ}{{\CYRHRDSN}}1 {Ы}{{\CYRY}}1
           {Ь}{{\CYRSFTSN}}1 {Э}{{\CYREREV}}1 {Ю}{{\CYRYU}}1 {Я}{{\CYRYA}}1
           {—}{{\textemdash}}1 {«}{{\guillemotleft}}1 {»}{{\guillemotright}}1,
}

% -- Theorem environments ---------------------------------------------
\newtheorem{definition}{Определение}[section]
\newtheorem{proposition}{Утверждение}[section]
\newtheorem{remark}{Замечание}[section]

% -- Custom commands --------------------------------------------------
\newcommand{\bs}{\mathrm{BS}}
\newcommand{\bss}{\mathrm{BSS}}
\DeclareMathOperator*{\argmin}{arg\,min}
\DeclareMathOperator*{\argmax}{arg\,max}

% =====================================================================
\title{%
  Delphi Press\\
  Методология прогнозирования заголовков СМИ%
}
\author{Delphi Press --- внутренний документ команды}
\date{Версия 0.9.5 \quad|\quad \today}

\begin{document}
\maketitle

\begin{abstract}
Документ описывает полную методологию системы Delphi Press ---
веб-продукта для прогнозирования заголовков СМИ.
Мультиагентный Дельфи-пайплайн (5 персон, 2 раунда, 18 агентов, 9 стадий)
обрабатывает данные информагенств и рынков предсказаний,
генерируя стилизованные заголовки для конкретных изданий.

Документ покрывает: архитектуру pipeline, метод Дельфи с 5 экспертными персонами,
сбор и сжатие данных, обогащение сигналами Polymarket (informed consensus, BSS $+0.196$),
LLM-инфраструктуру, систему оценки качества,
а также dead ends и уроки из 6 месяцев разработки.

Промпты всех агентов приведены в Приложении~B.
\end{abstract}

\tableofcontents
\newpage

% =====================================================================
% PARTS I + VII: System Overview + LLM Infrastructure
% =====================================================================
% ============================================================================
% PART I: OVERVIEW AND ARCHITECTURE
% ============================================================================

\section{Что такое Delphi Press}

\label{sec:overview}

% Source: 00-overview.md:1-10, CLAUDE.md:1-15

\emph{Delphi Press} — веб-продукт для автоматизированного прогнозирования заголовков средств массовой информации на заданную целевую дату. Система вводит пользователя в интерактивный мультиагентный пайплайн, который анализирует текущее состояние геополитических, экономических и медийных контекстов, а затем генерирует ранжированные предсказания заголовков с уровнями уверенности и обоснованиями.

\subsection{Проблемная постановка}

Исследователи-политологи и аналитики медиа сталкиваются с проблемой предсказания того, на какие события обратит внимание конкретное издание в конкретный день. Это требует:

\begin{itemize}[nosep]
\item Анализа текущих информационных потоков и планируемых событий
\item Понимания редакционной позиции и стилевых предпочтений издания
\item Моделирования вероятности события и его медийной значимости
\item Генерации заголовков, стилистически соответствующих целевому СМИ
\end{itemize}

Ручной анализ требует дней работы. \emph{Delphi Press} автоматизирует весь процесс через специализированный \emph{ensemble forecasting} подход — мультиагентный конвейер, вдохновлённый методом Дельфи.

\subsection{Два режима работы}

Система поддерживает две архитектуры развёртывания, приспособленные к разным пользовательским сценариям:

\begin{table}[h]
\centering
\begin{tabular}{|l|c|c|}
\hline
\textbf{Характеристика} & \textbf{Web UI} & \textbf{Claude Code} \\
\hline
Доступ & Веб-интерфейс & CLI (/predict в Claude Code) \\
API-ключ & Пользователь вводит свой (OpenRouter) & Покрыто подпиской Code Max \\
Стоимость прогноза & \$5–50 & \$0 \\
Модельное разнообразие & Нет (все Opus 4.6 (в коде: пресет \texttt{full}), diversity через промпты) & Нет (все Opus 4.6) \\
Автоматизация & REST API + расписание & Ручной запуск \\
Персистентность & БД, история, аналитика & Markdown-отчёт \\
\hline
\end{tabular}
\caption{Два режима Delphi Press. Оба разделяют промпты персон и Pydantic-схемы данных.}
\label{tab:modes}
\end{table}

\subsection{Целевая аудитория и язык}

Основная аудитория — исследователи-политологи, аналитики медиа, специалисты по молодёжной политике и стратегам госучреждений. Интерфейс разработан для нетехнических пользователей с интуитивной системой заполнения форм и визуализацией результатов.

Язык интерфейса — русский. Анализ проводится на медиа на любом языке; результаты генерируются на языке целевого СМИ.

\subsection{Технический стек}

\begin{table}[h]
\centering
\small
\begin{tabular}{|l|l|c|}
\hline
\textbf{Слой} & \textbf{Технология} & \textbf{Версия} \\
\hline
Язык & Python & 3.12+ \\
Backend & FastAPI + Uvicorn & 0.115+ \\
Task Queue & ARQ (Redis) & 0.26+ \\
База данных & SQLite + aiosqlite & — \\
ORM & SQLAlchemy & 2.0+ \\
Валидация & Pydantic & v2 \\
LLM клиент & OpenAI SDK (OpenRouter) & 1.0+ \\
HTTP клиент & httpx & 0.27+ \\
Frontend & Tailwind CSS + Vanilla JS & 4.2.2 \\
Шаблоны & Jinja2 & 3.1+ \\
Развёртывание & Docker Compose & 2.0+ \\
Шифрование & cryptography (Fernet) & 43+ \\
Авторизация & PyJWT + bcrypt & — \\
Тестирование & pytest + pytest-asyncio & — \\
\hline
\end{tabular}
\caption{Технический стек Delphi Press v0.9.5.}
\label{tab:tech_stack}
\end{table}

% Source: CLAUDE.md:10

\textbf{Версия}: 0.9.5. \textbf{Покрытие тестами}: 1324 тестов. \textbf{Принцип архитектуры}: модульный монолит — единый Python-проект с чёткими границами между модулями, общающимися через прямые вызовы функций (не HTTP RPC).

\subsection{Развёртывание: 4 контейнера Docker}

Система развёртывается как четыре скоординированных контейнера:

\begin{enumerate}[nosep]
\item \texttt{app} (FastAPI/Uvicorn на порту 8000) — веб-интерфейс, REST API, SSE-потоки
\item \texttt{worker} (ARQ на Redis) — асинхронное выполнение пайплайна прогнозирования
\item \texttt{redis} (порт 6379) — брокер задач, pub/sub для SSE
\item \texttt{nginx} (порты 80/443) — reverse proxy, SSL termination via Let's Encrypt
\end{enumerate}

\noindent Данные хранятся в SQLite-файле на диске (\texttt{/app/data/foresighting.db}). Архитектура выбрана для простоты развёртывания на одном VPS (Debian 12, 4 vCPU @ 20%, 8 GB RAM).


\section{Архитектура pipeline}

\label{sec:architecture}

% Source: architecture.md:1-50, 00-overview.md:266-400

\emph{Delphi Press} реализует девятистадийный \emph{sequential-parallel hybrid} конвейер: каждая стадия выполняется последовательно, но агенты внутри стадии могут работать параллельно (в зависимости от конфигурации). Контекст (состояние) передаётся через 16 типизированных слотов (\emph{PipelineContext}), заполняемых последовательно.

\subsection{Обзор стадий}

% Source: architecture.md:36-48

\begin{table}[h]
\centering
\small
\begin{tabular}{|c|l|l|c|c|c|}
\hline
\# & \textbf{Стадия} & \textbf{Агенты} & \textbf{Параллель} & \textbf{min\_successful} & \textbf{Timeout} \\
\hline
1 & COLLECTION & 4 агента & Да & 2/4 & 600s \\
2 & EVENT\_IDENTIFICATION & EventTrendAnalyzer & Нет & — & 600s \\
3 & TRAJECTORY & 3 аналитика & Да & 2/3 & 600s \\
4 & DELPHI\_R1 & 5 персон & Да & 3/5 & 600s \\
5 & DELPHI\_R2 & Медиатор+5 персон & Смешано & 3/5 & 900s \\
6 & CONSENSUS & Judge & Нет & — & 300s \\
7 & FRAMING & FramingAnalyzer & Нет & — & 300s \\
8 & GENERATION & StyleReplicator & Нет & — & 300s \\
9 & QUALITY\_GATE & QualityGate & Нет & — & 300s \\
\hline
\end{tabular}
\caption{Девять стадий пайплайна Delphi Press. Стадия 5 имеет двухфазную логику: сначала Медиатор (последовательно), затем 5 персон (параллельно, min=3).}
\label{tab:stages}
\end{table}

\subsection{Полная таблица агентов и LLM-задач}

% Source: architecture.md:11-30

\begin{table}[h]
\centering
\tiny
\begin{tabular}{|l|l|c|l|l|}
\hline
\textbf{Агент} & \textbf{Стадия} & \textbf{LLM-задачи} & \textbf{Слоты контекста} & \textbf{Файл} \\
\hline
NewsScout & 1 & news\_scout\_search & signals & collectors/news\_scout.py \\
EventCalendar & 1 & event\_calendar, event\_assessment & scheduled\_events & collectors/event\_calendar.py \\
OutletHistorian & 1 & outlet\_historian & outlet\_profile & collectors/outlet\_historian.py \\
ForesightCollector & 1 & (нет LLM) & foresight\_events, foresight\_signals & collectors/foresight\_collector.py \\
EventTrendAnalyzer & 2 & event\_clustering, trajectory\_analysis, cross\_impact\_analysis & event\_threads, trajectories, cross\_impact\_matrix & analysts/event\_trend.py \\
GeopoliticalAnalyst & 3 & geopolitical\_analysis & event\_threads[] & analysts/geopolitical.py \\
EconomicAnalyst & 3 & economic\_analysis & event\_threads[] & analysts/economic.py \\
MediaAnalyst & 3 & media\_analysis & event\_threads[] & analysts/media.py \\
DelphiRealist & 4/5 & delphi\_r1\_realist, delphi\_r2\_realist & round1/2\_assessments & forecasters/personas.py \\
DelphiGeostrateg & 4/5 & delphi\_r1\_geostrateg, delphi\_r2\_geostrateg & round1/2\_assessments & forecasters/personas.py \\
DelphiEconomist & 4/5 & delphi\_r1\_economist, delphi\_r2\_economist & round1/2\_assessments & forecasters/personas.py \\
DelphiMediaExpert & 4/5 & delphi\_r1\_media, delphi\_r2\_media & round1/2\_assessments & forecasters/personas.py \\
DelphiDevilsAdvocate & 4/5 & delphi\_r1\_devils, delphi\_r2\_devils & round1/2\_assessments & forecasters/personas.py \\
Mediator & 5 & mediator & mediator\_synthesis & forecasters/mediator.py \\
Judge & 6 & (детерминированный) & predicted\_timeline, ranked\_predictions & forecasters/judge.py \\
FramingAnalyzer & 7 & framing & framing\_briefs & generators/framing.py \\
StyleReplicator & 8 & style\_generation, style\_generation\_ru, style\_generation\_en & generated\_headlines & generators/style\_replicator.py \\
QualityGate & 9 & quality\_factcheck, quality\_style & final\_predictions & generators/quality\_gate.py \\
\hline
\end{tabular}
\caption{18 агентов, 28 LLM-задач, распределение по стадиям и слотам контекста.}
\label{tab:agents_full}
\end{table}

\subsection{Поток данных: примеры слотов}

% Source: architecture.md:116-139

Основные слоты \emph{PipelineContext}:

\begin{itemize}[nosep]
\item \textbf{signals}: List[SignalRecord] — сырые новости из RSS и поиска (100–200)
\item \textbf{scheduled\_events}: List[ScheduledEvent] — политические/экономические события на target\_date
\item \textbf{outlet\_profile}: OutletProfile — стилевой профиль издания с примерами заголовков
\item \textbf{event\_threads}: List[EventThread] — top-20 кластеризованных событийных цепочек
\item \textbf{trajectories}: List[EventTrajectory] — сценарные траектории (baseline, optimistic, pessimistic, black\_swan)
\item \textbf{cross\_impact\_matrix}: CrossImpactMatrix — матрица перекрёстных влияний между событиями
\item \textbf{round1\_assessments}: List[PersonaAssessment] — независимые оценки от 5 персон (раунд 1)
\item \textbf{mediator\_synthesis}: MediatorSynthesis — консенсус, расхождения, ключевые вопросы
\item \textbf{round2\_assessments}: List[PersonaAssessment] — пересмотренные оценки после медиации
\item \textbf{ranked\_predictions}: List[RankedPrediction] — top-7 финальных событий, ранжированных по headline\_score
\item \textbf{framing\_briefs}: List[FramingBrief] — анализ подачи каждого события
\item \textbf{generated\_headlines}: List[GeneratedHeadline] — сгенерированные заголовки и абзацы (2–3 на событие)
\item \textbf{final\_predictions}: List[FinalPrediction] — финальные прогнозы после факт-чека и дедупликации
\end{itemize}

\subsection{Прямолинейный пример потока}

\begin{equation}
\text{PredictionRequest} \xrightarrow{\text{Stage 1}} \text{signals, events, outlet\_profile}
\end{equation}

\begin{equation}
\xrightarrow{\text{Stage 2}} \text{event\_threads, trajectories, cross\_impact\_matrix}
\end{equation}

\begin{equation}
\xrightarrow{\text{Stages 3–5}} \text{round1/2\_assessments, mediator\_synthesis}
\end{equation}

\begin{equation}
\xrightarrow{\text{Stages 6–9}} \text{ranked\_predictions} \to \text{framing\_briefs} \to \text{generated\_headlines} \to \text{final\_predictions}
\end{equation}

\begin{equation}
\xrightarrow{\text{Response}} \text{PredictionResponse} \text{ с 7 заголовков, уверенностью, обоснованиями}
\end{equation}

\subsection{Производительность и стоимость}

% Source: 00-overview.md:435-450, CLAUDE.md:10-11

\begin{table}[h]
\centering
\begin{tabular}{|l|c|c|c|}
\hline
\textbf{Стадия} & \textbf{LLM-вызовов} & \textbf{Основная модель} & \textbf{Стоимость} \\
\hline
1: Collection & ~5 & Gemini + Sonnet & \$1.50 \\
2: Event Identification & ~25 & Gemini-flash & \$0.50 \\
3: Trajectory & ~21 & Sonnet & \$3.00 \\
4: Delphi R1 & 5 & Mix (5 моделей) & \$8.00 \\
5a: Mediator & 1 & Opus & \$2.00 \\
5b: Delphi R2 & 5 & Mix (5 моделей) & \$8.00 \\
6: Judge & 1 & Opus & \$2.00 \\
7: Framing & ~7 & Sonnet & \$2.00 \\
8: Generation & ~21 & Sonnet & \$1.50 \\
9: Quality Gate & ~14 & Sonnet & \$2.00 \\
\hline
\textbf{Итого} & \textbf{~105 вызовов} & & \textbf{~\$30.50} \\
\hline
\end{tabular}
\caption{Бюджет LLM на один прогноз. Время выполнения: 15–20 минут на одном VPS.}
\label{tab:cost_budget}
\end{table}


\section{Принципы проектирования}

\label{sec:design_principles}

% Source: 00-overview.md:26-35, 02-agents-core.md:10-16, CLAUDE.md:72-81

\subsection{Модульный монолит, а не микросервисы}

Система построена как \emph{unified codebase} с чёткой модульной структурой. Каждый модуль:

\begin{itemize}[nosep]
\item Живёт в отдельной директории (\texttt{src/agents/collectors/}, \texttt{src/agents/analysts/} и т. д.)
\item Имеет определённый контракт: типизированный вход (Pydantic-модель), выход (\texttt{AgentResult})
\item Реализуется и тестируется независимо
\item Общается с другими модулями через прямые вызовы функций (внутри одного процесса)
\end{itemize}

\noindent Этот подход избегает сложности микросервисной архитектуры (сетевые задержки, согласованность состояния) при сохранении модульности и тестируемости.

\subsection{\emph{Async everywhere} — асинхронная обработка I/O}

% Source: CLAUDE.md:74

Все операции ввода-вывода реализованы асинхронно:

\begin{itemize}[nosep]
\item \textbf{httpx} — асинхронный HTTP-клиент для RSS, web search, скрейпинга
\item \textbf{aiosqlite} — асинхронный драйвер SQLite для БД-операций
\item \textbf{ARQ} — асинхронная очередь задач на Redis
\item \textbf{asyncio.gather()} — параллельное выполнение агентов внутри стадии
\end{itemize}

\noindent Вся вычислительная база от app-сервера до worker-процесса построена на \texttt{asyncio}, что позволяет эффективно обрабатывать множество параллельных прогнозов на одном VPS.

\subsection{\emph{Fail-soft} с \texttt{min\_successful} логикой}

Система не требует, чтобы \textit{все} агенты в стадии завершились успешно. Вместо этого определяется минимальное количество успешных агентов:

\begin{itemize}[nosep]
\item Stage 1 (Collection): min=2/4 агентов — система может продолжить даже если 2 коллектора отказали
\item Stage 3 (Trajectory): min=2/3 аналитиков — достаточно 2 аспектов анализа
\item Stage 4 (Delphi R1): min=3/5 персон — мнение 3 экспертов из 5 достаточно для устойчивого консенсуса
\item Stage 5 (Delphi R2): min=3/5 (после медиации)
\end{itemize}

\noindent Это обеспечивает \emph{graceful degradation}: если primary-модель недоступна, \texttt{ModelRouter} переключается на fallback-модель через OpenRouter и продолжает работу.

\subsection{Разнообразие через промпты и когнитивные смещения}

% Source: router.py:224 "Diversity обеспечивается промптами и когнитивными смещениями, не моделями."

В текущей реализации все 5 персон Дельфи используют \textbf{одну и ту же модель} (Claude Opus 4.6 через OpenRouter). Разнообразие ошибок обеспечивается не выбором разных LLM, а тремя механизмами:

\begin{enumerate}[nosep]
\item \textbf{Уникальные системные промпты} ($\sim$400--500 строк на персону) --- каждая персона получает свой аналитический фреймворк, фокус и стиль рассуждений
\item \textbf{Когнитивные смещения} --- для каждой персоны определены \texttt{over\_predicts}, \texttt{under\_predicts} и \texttt{anchor\_type} (см. Часть~III)
\item \textbf{Горизонт-адаптивные веса} --- Judge корректирует влияние каждой персоны в зависимости от горизонта прогнозирования (см. Часть~V)
\end{enumerate}

\noindent Модельное разнообразие (5 разных LLM) рассматривалось как альтернатива, но не реализовано --- см. Часть~IX (Dead Ends) для обоснования.

\subsection{Pydantic v2 для валидации и структурированного вывода}

% Source: CLAUDE.md:76-77

Все входы и выходы агентов типизированы через Pydantic-модели (\texttt{src/schemas/*}):

\begin{itemize}[nosep]
\item \textbf{PredictionRequest} --- ввод пользователя (выход издания, целевая дата)
\item \textbf{AgentResult} $\to$ выход каждого агента (данные, метрики, ошибки)
\item \textbf{PipelineContext} $\to$ разделяемое состояние (16 слотов)
\item \textbf{PersonaAssessment}, \textbf{RankedPrediction}, \textbf{FinalPrediction} $\to$ типизированные структуры результатов
\end{itemize}

\noindent \textbf{Structured Output}: LLM-ответы не парсятся как сырые строки. Вместо этого OpenRouter запрашивает JSON output, который затем валидируется через Pydantic с помощью \texttt{json\_mode=true}. Это снижает вероятность парсинга неверного формата.

\subsection{Стоимость и бюджетный контроль}

% Source: CLAUDE.md:11, architecture.md:236-239

Каждый вызов LLM отслеживает:

\begin{itemize}[nosep]
\item Использованные токены (входные/выходные)
\item Стоимость в USD (по таблице прайсинга OpenRouter)
\item Модель и провайдер
\item Время выполнения
\end{itemize}

\noindent Глобальный лимит: \textbf{\$50 USD на один прогноз} (настраивается). Если BudgetTracker обнаруживает, что предстоящий вызов превысит лимит, система выбрасывает \texttt{LLMBudgetExceededError} и прерывает пайплайн.


% ============================================================================
% PART VII: LLM INFRASTRUCTURE
% ============================================================================

\section{Интеграция с OpenRouter}

\label{sec:openrouter}

% Source: 07-llm-layer.md:1-30, architecture.md:230-231, CLAUDE.md:11

\emph{Delphi Press} использует \emph{OpenRouter.ai} — унифицированный OpenAI-совместимый API для доступа к 200+ LLM-моделям (Claude, GPT-4, Gemini, Llama, DeepSeek и др.) через единый эндпоинт.

\subsection{Архитектура провайдера}

\begin{itemize}[nosep]
\item \textbf{Base URL}: \texttt{https://openrouter.ai/api/v1}
\item \textbf{Формат моделей}: \texttt{provider/model}, например \texttt{anthropic/claude-sonnet-4}, \texttt{openai/gpt-4o}
\item \textbf{Клиент}: OpenAI Python SDK v1.0+ с переопределённым \texttt{base\_url}
\item \textbf{Аутентификация}: API-ключ в заголовке \texttt{Authorization: Bearer sk-or-...}
\end{itemize}

\subsection{Таблица моделей и задач}

% Source: architecture.md:54-114, 07-llm-layer.md:458-478

Все 28 LLM-задач пайплайна настроены с primary и fallback моделями:

\begin{table}[h]
\centering
\tiny
\begin{tabular}{|l|l|l|l|c|c|}
\hline
\textbf{Task ID} & \textbf{Назначение} & \textbf{Primary Model} & \textbf{Fallback} & \textbf{Temp} & \textbf{JSON} \\
\hline
\multicolumn{6}{|c|}{\textbf{COLLECTORS (Stage 1)}} \\
\hline
news\_scout\_search & Формирование поисковых запросов & Gemini 3.1-flash-lite & Gemini 2.5-flash & 0.3 & Net \\
event\_calendar & Поиск запланированных событий & Gemini 3.1-flash-lite & Gemini 2.5-flash & 0.3 & Da \\
event\_assessment & Оценка значимости событий & Claude Opus 4.6 & Claude Sonnet 4.5 & 0.4 & Da \\
outlet\_historian & Анализ стиля издания & Claude Opus 4.6 & Claude Sonnet 4.5 & 0.4 & Da \\
\hline
\multicolumn{6}{|c|}{\textbf{ANALYSTS (Stages 2–3)}} \\
\hline
event\_clustering & Кластеризация сигналов & Gemini 3.1-flash-lite & Gemini 2.5-flash & 0.2 & Da \\
trajectory\_analysis & Сценарии развития событий & Claude Opus 4.6 & Claude Sonnet 4.5 & 0.6 & Da \\
cross\_impact\_analysis & Матрица перекрёстных влияний & Claude Opus 4.6 & Claude Sonnet 4.5 & 0.4 & Da \\
geopolitical\_analysis & Геополитический контекст & Claude Opus 4.6 & Claude Sonnet 4.5 & 0.5 & Da \\
economic\_analysis & Экономический контекст & Claude Opus 4.6 & Claude Sonnet 4.5 & 0.5 & Da \\
media\_analysis & Медийный контекст & Claude Opus 4.6 & Claude Sonnet 4.5 & 0.5 & Da \\
\hline
\multicolumn{6}{|c|}{\textbf{DELPHI R1 (Stage 4)}} \\
\hline
delphi\_r1\_realist & Реалист & Claude Opus 4.6 & Claude Sonnet 4.5 & 0.7 & Da \\
delphi\_r1\_geostrateg & Геостратег & Claude Opus 4.6 & Claude Sonnet 4.5 & 0.7 & Da \\
delphi\_r1\_economist & Экономист & Claude Opus 4.6 & Claude Sonnet 4.5 & 0.7 & Da \\
delphi\_r1\_media & Медиа-эксперт & Claude Opus 4.6 & Claude Sonnet 4.5 & 0.7 & Da \\
delphi\_r1\_devils & Адвокат дьявола & Claude Opus 4.6 & Claude Sonnet 4.5 & 0.9 & Da \\
\hline
\multicolumn{6}{|c|}{\textbf{MEDIATOR \& DELPHI R2 (Stage 5)}} \\
\hline
mediator & Синтез расхождений & Claude Opus 4.6 & Claude Sonnet 4.5 & 0.5 & Da \\
delphi\_r2\_* & Дельфи раунд 2 (пересмотр) & Claude Opus 4.6 & Claude Sonnet 4.5 & 0.6 & Da \\
\hline
\multicolumn{6}{|c|}{\textbf{GENERATORS (Stages 6–9)}} \\
\hline
framing & Анализ фрейминга & Claude Opus 4.6 & Claude Sonnet 4.5 & 0.5 & Da \\
style\_generation & Генерация заголовков (мультиязык) & Claude Opus 4.6 & Claude Sonnet 4.5 & 0.8 & Net \\
style\_generation\_ru & Генерация заголовков (русский) & Claude Opus 4.6 & Claude Sonnet 4.5 & 0.8 & Net \\
style\_generation\_en & Генерация заголовков (английский) & Claude Opus 4.6 & Claude Sonnet 4.5 & 0.8 & Net \\
quality\_factcheck & Факт-чек & Claude Opus 4.6 & Claude Sonnet 4.5 & 0.2 & Da \\
quality\_style & Стилистическая проверка & Claude Opus 4.6 & Claude Sonnet 4.5 & 0.3 & Da \\
\hline
\end{tabular}
\caption{28 LLM-задач с моделями и параметрами. JSON mode включён для структурированного вывода (схемы Pydantic).}
\label{tab:task_models}
\end{table}

\subsection{JSON Mode для структурированного вывода}

% Source: 07-llm-layer.md:90-104

Большинство задач (24 из 28) используют JSON mode. Это означает:

\begin{itemize}[nosep]
\item OpenRouter просит модель вывести JSON
\item Ответ парсится как \texttt{dict} и валидируется через Pydantic-модель
\item Если парсинг или валидация не удаётся, агент возвращает \texttt{AgentResult(success=False)}
\end{itemize}

\noindent Это обеспечивает \emph{deterministic parsing} и снижает риск галлюцинаций в неправильном формате.



\section{Retry, стоимость и бюджетный контроль}

\label{sec:retry_budget}

% Source: 07-llm-layer.md:366-445, 649-707

\subsection{\emph{Exponential backoff} с jitter}

% Source: 07-llm-layer.md:366-408

При ошибке LLM-провайдера (HTTP 429, 500, 502, 503, 504) система выполняет повторный запрос с экспоненциальной задержкой:

\begin{equation}
\text{delay}_n = \min\left(\text{base} \cdot 2^n, \text{max\_delay}\right) + \text{jitter}
\end{equation}

где:
\begin{itemize}[nosep]
\item base = 1 сек (базовая задержка)
\item max\_delay = 30 сек (потолок)
\item jitter $\in$ [0, base/2] = [0, 0.5] сек (случайный разброс для избежания thundering herd)
\item n = номер попытки (0, 1, 2, ...)
\end{itemize}

Например:
\begin{itemize}[nosep]
\item Попытка 0: сразу выполнить
\item Попытка 1: ошибка $\to$ ждём $1 \cdot 2^0$ + jitter = 1--1.5 сек
\item Попытка 2: ошибка $\to$ ждём $1 \cdot 2^1$ + jitter = 2--2.5 сек
\item Попытка 3: ошибка $\to$ ждём $1 \cdot 2^2$ + jitter = 4--4.5 сек
\item Попытка 4 (последняя): ошибка $\to$ всё, выбрасываем исключение
\end{itemize}

\subsection{Fallback-цепочка}

% Source: 07-llm-layer.md:622-645

При исчерпании всех retry-попыток для primary модели система переключается на fallback-модели по очереди:

\begin{equation}
\begin{aligned}
&\text{Primary Model}\\
&\quad \xrightarrow{\text{retry $\times$ max\_retries}} \text{Fallback[0]}\\
&\quad \xrightarrow{\text{retry $\times$ max\_retries}} \text{Fallback[1]}\\
&\quad \xrightarrow{\text{...}} \text{Fallback[n]}\\
&\quad \xrightarrow{\text{all fail}} \text{raise LLMProviderError}
\end{aligned}
\end{equation}

Пример для \texttt{delphi\_r1\_realist}:
\begin{enumerate}[nosep]
\item Попробовать Claude Opus 4.6 (primary) $\times$ 3 retry
\item Если всё ещё ошибка $\to$ попробовать Claude Sonnet 4.5 (fallback) $\times$ 3 retry
\item Если оба отказали $\to$ выбросить исключение, агент вернёт \texttt{AgentResult(success=False)}
\end{enumerate}

\subsection{HTTP 429 (Rate Limiting)}

Если провайдер возвращает HTTP 429 с заголовком \texttt{Retry-After}, система уважает этот заголовок:

\begin{equation}
\text{delay} = \max(\text{Retry-After}, \text{exponential backoff})
\end{equation}

\subsection{Формула стоимости вызова}

% Source: 07-llm-layer.md:316-327, CLAUDE.md:11

Каждый LLM-вызов рассчитывает стоимость по формуле:

\begin{equation}
\text{cost\_usd} = \left(\frac{\text{tokens\_in}}{1{,}000{,}000} \times \text{price\_in}\right) + \left(\frac{\text{tokens\_out}}{1{,}000{,}000} \times \text{price\_out}\right)
\label{eq:llm_cost}
\end{equation}

где:
\begin{itemize}[nosep]
\item \texttt{tokens\_in}, \texttt{tokens\_out} — переданные OpenRouter в \texttt{usage}
\item \texttt{price\_in}, \texttt{price\_out} --- цены из таблицы \texttt{MODEL\_PRICING} (\$/1M токенов)
\end{itemize}

Пример: Claude Opus 4.6
\begin{itemize}[nosep]
\item Input: \$3 за 1М токенов
\item Output: \$15 за 1М токенов
\item Для 10k входных, 1k выходных: cost = (10k/1M $\times$ 3) + (1k/1M $\times$ 15) = 0.03 + 0.015 = \$0.045
\end{itemize}

\subsection{Бюджетный контроль}

% Source: 07-llm-layer.md:649-707, architecture.md:236-239

\emph{BudgetTracker} отслеживает расходы в рамках одного прогноза и блокирует вызовы, которые превысят лимит.

\begin{table}[h]
\centering
\begin{tabular}{|l|l|}
\hline
\textbf{Параметр} & \textbf{Значение} \\
\hline
Лимит на прогноз & \$50 USD (по умолчанию, настраивается) \\
Проверка перед вызовом & Если \texttt{estimated\_cost > remaining\_budget}, raise \texttt{LLMBudgetExceededError} \\
Типичный прогноз & \$5–15 USD \\
Максимальный прогноз & До \$50 USD (при многих событиях, долгих раундах) \\
\hline
\end{tabular}
\caption{Бюджетный контроль на уровне провайдера.}
\label{tab:budget}
\end{table}

Алгоритм:

\begin{equation}
\text{remaining} = \text{budget\_usd} - \sum_{i=1}^{n} \text{cost\_usd}_i
\end{equation}

\begin{equation}
\text{if } \text{estimated\_cost} > \text{remaining} \text{ then } \text{raise BudgetExceededError}
\end{equation}

Оценка стоимости рассчитывается по количеству токенов в промпте:

\begin{equation}
\text{estimated\_tokens\_in} = \text{estimate\_messages\_tokens}(\text{messages})
\end{equation}

\begin{equation}
\text{estimated\_cost} = \frac{\text{estimated\_tokens\_in}}{1{,}000{,}000} \times \text{price\_in} \times 1.2
\end{equation}

\noindent Коэффициент 1.2 добавляет буфер на выходные токены (обычно выход меньше входа, но разница варьируется).

\subsection{Примеры и мониторинг}

\emph{BudgetTracker} предоставляет методы для аналитики:

\begin{itemize}[nosep]
\item \texttt{summary\_by\_stage()}: расходы по стадиям пайплайна
\item \texttt{summary\_by\_model()}: расходы по моделям
\item \texttt{to\_records()}: все \texttt{CostRecord} для сохранения в БД
\end{itemize}

Пример логирования:

\begin{verbatim}
Стадия 1 (Collection): $1.50
Стадия 2 (Event Identification): $0.50
Стадия 3 (Trajectory): $3.00
Стадия 4 (Delphi R1): $8.00
Стадия 5 (Mediator + R2): $10.00
Стадия 6 (Consensus): $2.00
Стадия 7–9 (Framing + Generation + QG): $5.50
-----------
Итого: $30.50 USD
\end{verbatim}


\section{Список аббревиатур}

\begin{table}[h]
\centering
\small
\begin{tabular}{|l|l|}
\hline
\textbf{Аббревиатура} & \textbf{Расшифровка} \\
\hline
API & Application Programming Interface \\
ARQ & Async Redis Queue \\
CSS & Cascading Style Sheets \\
Delphi & Метод экспертного консенсуса (делфийский метод) \\
HTTP & HyperText Transfer Protocol \\
JSON & JavaScript Object Notation \\
JWT & JSON Web Token \\
LLM & Large Language Model \\
ORM & Object-Relational Mapping \\
REST & Representational State Transfer \\
RSS & Really Simple Syndication \\
SQLite & SQL Lite (встроённая база данных) \\
SSE & Server-Sent Events \\
UUID & Universally Unique Identifier \\
VPS & Virtual Private Server \\
\hline
\end{tabular}
\caption{Ключевые аббревиатуры, используемые в тексте.}
\label{tab:abbrev}
\end{table}

% End of Parts I + VII


\newpage
% =====================================================================
% PART II: Data Collection (Stages 1-2)
% =====================================================================
% PART II: Сбор данных (Stages 1-2)
% Source: docs/03-collectors.md, docs/01-data-sources.md, docs/04-analysts.md

\section{NewsScout: агрегация сигналов}

Первый этап сбора информации в пайплайне Delphi Press основан на агенте~\texttt{NewsScout}, чья задача~--- собрать многоканальные сигналы из различных источников и сформировать унифицированный набор записей для последующей обработки.

% Source: 03-collectors.md:477-489

\subsection{Архитектура сбора}

\texttt{NewsScout} реализует двухисточниковую стратегию:

\begin{enumerate}
    \item \textbf{RSS-агрегация}: параллельная загрузка 20--30 RSS-фидов из мировых, региональных и тематических источников
    \begin{itemize}
        \item Глобальные источники (Reuters, AP, BBC, Al Jazeera и др.) --- 10--15 фидов
        \item Региональные (ТАСС, РИА для русскоязычных; CNN, NYT для англоязычных) --- 5--10
        \item Тематические (финансы, политика, наука) --- 5--10
    \end{itemize}

    \item \textbf{Веб-поиск}: параллельное выполнение 3--5 поисковых запросов через API Exa/Jina MCP
    \begin{itemize}
        \item Общий запрос: \emph{major news events {target\_date} world}
        \item Региональный: \emph{news {region\_of\_outlet} {target\_date}}
        \item Тематический: \emph{scheduled events politics economy {target\_date}}
        \item Outlet-specific: \emph{topics covered by {outlet} this week}
        \item Трендовый: \emph{breaking news today {current\_date}} (для контекста)
    \end{itemize}
\end{enumerate}

Оба источника обрабатываются параллельно с помощью \texttt{asyncio.gather()}. Если один источник недоступен (например, RSS-фид не отвечает в течение 30 секунд), пайплайн продолжает работу с другим, логируя ошибку. Критическое условие: если оба источника полностью недоступны, возвращается исключение.

% Source: 03-collectors.md:542-601

\subsection{Дедупликация и релевантность}

После объединения результатов из обоих источников выполняется дедупликация по URL:

\begin{itemize}
    \item Для каждого сигнала вычисляется нормализованный URL (строчные буквы, удаление trailing slash)
    \item При обнаружении дубля сохраняется сигнал с более высоким \texttt{relevance\_score}
\end{itemize}

Релевантность сигнала рассчитывается на основе:
\begin{itemize}
    \item Для RSS-сигналов: свежесть (дата публикации) и надёжность источника
    \item Для веб-поиска: \texttt{score} из API поисковика (нормализуется в диапазон 0--1)
\end{itemize}

Опциональная LLM-классификация активируется для сигналов без категорий. Используется модель \texttt{openai/gpt-4o-mini} с пакетной обработкой (batch по 20 сигналов). Промпт требует классифицировать заголовки по доменам (политика, экономика, наука и т.д.) и извлечь именованные сущности (названия людей, организаций, мест).

% Source: 03-collectors.md:53-129

\subsection{Схема SignalRecord}

Выход \texttt{NewsScout} --- список объектов \texttt{SignalRecord} (100--200 записей), определённых в \texttt{src/schemas/events.py}. Каждый сигнал содержит:

\begin{table}[h!]
\centering
\caption{Поля SignalRecord}
\label{tab:signal_record}
\small
\begin{tabular}{|l|l|l|}
\hline
\textbf{Поле} & \textbf{Тип} & \textbf{Назначение} \\
\hline
\texttt{id} & str & Уникальный ID: \texttt{sig\_\{hash8\}} \\
\texttt{title} & str & Заголовок (макс. 500 символов) \\
\texttt{summary} & str & Краткое содержание (макс. 1000 символов) \\
\texttt{url} & str & URL источника \\
\texttt{source\_name} & str & Название источника (Reuters, ТАСС, BBC и т.д.) \\
\texttt{source\_type} & Enum & RSS, WEB\_SEARCH, SOCIAL, WIRE \\
\texttt{published\_at} & datetime | None & Дата/время публикации (UTC) \\
\texttt{language} & str & ISO 639-1 код (ru, en, zh и т.д.) \\
\texttt{categories} & list[str] & Теги/категории (politics, economy, science и т.д.) \\
\texttt{entities} & list[str] & Именованные сущности (Trump, ЕЦБ, НАТО) \\
\texttt{relevance\_score} & float & Оценка релевантности (0.0--1.0) \\
\hline
\end{tabular}
\end{table}

% Source: 03-collectors.md:744-756

ID сигнала генерируется как SHA-256 хеш конкатенации URL и заголовка, обрезанный до первых 8 символов. Это обеспечивает детерминированность: один и тот же сигнал, встреченный дважды, получит одинаковый ID.

% Source: 03-collectors.md:533-536, 594-599

Завершающий этап --- сортировка по \texttt{relevance\_score} (убывающий порядок) и ограничение на 200 записей. Сигналы отсортированы от наиболее релевантных к наименее релевантным.

---

\section{EventCalendar: структурированные события}

Второй источник информации в Stage 1 --- агент \texttt{EventCalendar}, который выявляет запланированные события на целевую дату пруи с использованием веб-поиска и структурирования через LLM.

% Source: 03-collectors.md:787-830

\subsection{Логика поиска событий}

Процесс разбивается на пять этапов:

\begin{enumerate}
    \item \textbf{Формирование поисковых запросов} (5--8 запросов) по типам событий:
    \begin{itemize}
        \item Политика: \emph{political events scheduled {date}}, \emph{parliamentary sessions votes {date}}
        \item Экономика: \emph{economic data releases {date}}, \emph{central bank meetings {date}}, \emph{earnings reports {date}}
        \item Дипломатия: \emph{international summits meetings {date}}, \emph{UN sessions {date}}
        \item Суды: \emph{major court hearings verdicts {date}}
        \item Культура/спорт: \emph{major events conferences {date}}
    \end{itemize}
    Для русскоязычных СМИ часть запросов дублируется на русском.

    \item \textbf{Параллельный веб-поиск}: выполнение всех запросов асинхронно через Exa/Jina API

    \item \textbf{LLM-структурирование} (модель \texttt{openai/gpt-4o-mini}): извлечение из сырых результатов структурированных событий с полями title, description, event\_type, certainty, location, participants, potential\_impact

    \item \textbf{Дедупликация}: fuzzy matching по заголовкам (ratio > 0.8 по Levenshtein), совпадение типа и даты события

    \item \textbf{Оценка новостной значимости} (модель \texttt{anthropic/claude-sonnet-4}): ранжирование по \texttt{newsworthiness} с учётом вероятности освещения в выбранном издании и исключительности события
\end{enumerate}

% Source: 03-collectors.md:148-228

\subsection{Типология и уверенность}

Типы событий охватывают основные домены публичной жизни:

\begin{equation}
\text{EventType} \in \{\text{POLITICAL, ECONOMIC, DIPLOMATIC, JUDICIAL, MILITARY, CULTURAL, SCIENTIFIC, SPORTS, OTHER}\}
\end{equation}

Каждое событие сопровождается оценкой уверенности, которая отражает степень надёжности информации:

\begin{table}[h!]
\centering
\caption{Уровни уверенности в наступлении события}
\label{tab:event_certainty}
\begin{tabular}{|l|l|}
\hline
\textbf{Уровень} & \textbf{Интерпретация} \\
\hline
CONFIRMED & Официально подтверждено (расписание на сайте, пресс-релиз) \\
LIKELY & Высокая вероятность (по расписанию, анонсировано) \\
POSSIBLE & Возможно (слухи, предварительные анонсы) \\
SPECULATIVE & Спекулятивно (восстановлено из паттернов: ежемесячные отчёты) \\
\hline
\end{tabular}
\end{table}

% Source: 03-collectors.md:156-228

\subsection{Схема ScheduledEvent}

Выход \texttt{EventCalendar} --- список объектов \texttt{ScheduledEvent} (3--10 записей на типичный день, до 50 в день повышенной активности):

\begin{table}[h!]
\centering
\caption{Поля ScheduledEvent}
\label{tab:scheduled_event}
\small
\begin{tabular}{|l|l|l|}
\hline
\textbf{Поле} & \textbf{Тип} & \textbf{Назначение} \\
\hline
\texttt{id} & str & Уникальный ID: \texttt{evt\_\{hash\}} \\
\texttt{title} & str & Название события (макс. 300 символов) \\
\texttt{description} & str & Подробное описание (макс. 500 символов) \\
\texttt{event\_date} & date & Дата события (совпадает с target\_date) \\
\texttt{event\_type} & EventType & Тип события (political, economic и т.д.) \\
\texttt{certainty} & EventCertainty & Уровень уверенности (confirmed/likely/possible/speculative) \\
\texttt{location} & str & Место проведения \\
\texttt{participants} & list[str] & Ключевые участники/организации \\
\texttt{potential\_impact} & str & Потенциальное влияние на новостную повестку \\
\texttt{source\_url} & str & URL источника информации \\
\texttt{newsworthiness} & float & Оценка новостной значимости (0.0--1.0) \\
\hline
\end{tabular}
\end{table}

Значение \texttt{newsworthiness} = 1.0 означает, что событие гарантированно попадёт в топе новостей выбранного издания.

---

\section{OutletHistorian: профиль издания}

Третий источник данных в Stage 1 --- агент \texttt{OutletHistorian}, который анализирует стилевые характеристики и редакционную позицию целевого СМИ на основе исторических публикаций.

% Source: 03-collectors.md:1021-1076

\subsection{Методология анализа}

Анализ разбивается на три параллельных LLM-вызова:

\begin{enumerate}
    \item \textbf{Анализ стиля заголовков} (HeadlineStyle) --- 30 последних заголовков издания
    \item \textbf{Анализ стиля текстов} (WritingStyle) --- 10 первых абзацев статей
    \item \textbf{Анализ редакционной позиции} (EditorialPosition) --- полный контекст 50--100 статей за 30 дней
\end{enumerate}

Все три вызова выполняются одновременно через \texttt{asyncio.gather()}, минимизируя общее время выполнения. Модель используется единая: \texttt{anthropic/claude-sonnet-4} (для высокой точности лингвистического анализа).

Кеширование: если профиль для издания был создан менее 7 дней назад, он возвращается из кеша без переанализа. Это обосновано тем, что стилевые характеристики СМИ меняются медленно.

% Source: 03-collectors.md:323-383

\subsection{HeadlineStyle: метрики заголовков}

Характеристики стиля заголовков с точки зрения формы и содержания:

\begin{table}[h!]
\centering
\caption{Показатели HeadlineStyle}
\label{tab:headline_style}
\footnotesize
\begin{tabular}{|l|l|l|}
\hline
\textbf{Метрика} & \textbf{Тип} & \textbf{Описание} \\
\hline
\texttt{avg\_length\_chars} & int & Средняя длина заголовка (символы) \\
\texttt{avg\_length\_words} & int & Средняя длина заголовка (слова) \\
\texttt{uses\_colons} & bool & Двоеточие для разделения темы и деталя \\
\texttt{uses\_quotes} & bool & Цитаты в заголовках \\
\texttt{uses\_questions} & bool & Вопросительные заголовки \\
\texttt{uses\_numbers} & bool & Частое использование цифр (5 причин...) \\
\texttt{capitalization} & str & Стиль: sentence\_case, title\_case, all\_caps\_first\_word, lowercase \\
\texttt{vocabulary\_register} & str & Регистр: formal, neutral, colloquial, technical, mixed \\
\texttt{emotional\_tone} & str & Тон: neutral, alarming, optimistic, dramatic, ironic, dry \\
\texttt{common\_patterns} & list[str] & 3+ повторяющихся паттерна (например, \{Name\}: statement...) \\
\hline
\end{tabular}
\end{table}

% Source: 03-collectors.md:386-422

\subsection{WritingStyle: характеристики текстов}

Стилевые характеристики основного содержания статей:

\begin{table}[h!]
\centering
\caption{Показатели WritingStyle}
\label{tab:writing_style}
\footnotesize
\begin{tabular}{|l|l|l|}
\hline
\textbf{Поле} & \textbf{Тип} & \textbf{Описание} \\
\hline
\texttt{first\_paragraph\_style} & str & inverted\_pyramid, narrative, analytical, quote\_lead \\
\texttt{avg\_first\_paragraph\_sentences} & int & Среднее число предложений в первом абзаце \\
\texttt{avg\_first\_paragraph\_words} & int & Среднее число слов \\
\texttt{attribution\_style} & str & source\_first, source\_last, inline \\
\texttt{uses\_dateline} & bool & Наличие датлайна (МОСКВА, 2 апреля ---) \\
\texttt{paragraph\_length} & str & short, medium, long \\
\hline
\end{tabular}
\end{table}

\textit{Inverted pyramid} (перевёрнутая пирамида) --- классический журналистский стиль, когда самая важная информация (кто-что-где-когда) размещена в первом абзаце.

% Source: 03-collectors.md:425-472

\subsection{EditorialPosition: редакционная позиция}

Редакционная позиция и предпочтения издания определяют, какие темы получают приоритет и как они подаются:

\begin{table}[h!]
\centering
\caption{Компоненты EditorialPosition}
\label{tab:editorial_position}
\footnotesize
\begin{tabular}{|l|l|l|}
\hline
\textbf{Поле} & \textbf{Тип} & \textbf{Содержание} \\
\hline
\texttt{tone} & ToneProfile & neutral, conservative, liberal, sensationalist, analytical, official, oppositional \\
\texttt{focus\_topics} & list[str] & Приоритетные темы (внутренняя политика, экономика РФ и т.д.) \\
\texttt{avoided\_topics} & list[str] & Обходимые стороной темы \\
\texttt{framing\_tendencies} & list[str] & Типичные фреймы: pro\_government, market\_oriented, human\_interest, conflict\_frame \\
\texttt{source\_preferences} & list[str] & Предпочтительные источники (официальные лица, МИД, анонимные) \\
\texttt{stance\_on\_current\_topics} & dict & Позиция по ключевым темам (тема: позиция) \\
\texttt{omissions} & list[str] & Систематически пропускаемые темы \\
\hline
\end{tabular}
\end{table}

% Source: 03-collectors.md:289-320

\subsection{OutletProfile: итоговая структура}

Полный профиль издания объединяет все три анализа:

\begin{table}[h!]
\centering
\caption{Структура OutletProfile}
\label{tab:outlet_profile}
\footnotesize
\begin{tabular}{|l|l|l|}
\hline
\textbf{Поле} & \textbf{Тип} & \textbf{Назначение} \\
\hline
\texttt{outlet\_name} & str & Каноническое название издания \\
\texttt{outlet\_url} & str & Основной URL издания \\
\texttt{language} & str & ISO 639-1 код (ru, en и т.д.) \\
\texttt{headline\_style} & HeadlineStyle & Объект с метриками заголовков \\
\texttt{writing\_style} & WritingStyle & Объект с метриками текстов \\
\texttt{editorial\_position} & EditorialPosition & Объект с редакционной позицией \\
\texttt{sample\_headlines} & list[str] & 10--30 последних заголовков \\
\texttt{sample\_first\_paragraphs} & list[str] & 5--10 примеров первых абзацев \\
\texttt{analysis\_period\_days} & int & 30 (период анализа) \\
\texttt{articles\_analyzed} & int & Количество проанализированных статей \\
\texttt{analyzed\_at} & datetime & Timestamp проведения анализа \\
\hline
\end{tabular}
\end{table}

Профиль кешируется в SQLite таблице \texttt{outlet\_profiles} с ключом \texttt{outlet\_name} (нормализованное название) и TTL 7 дней.

---

\section{ForesightCollector: рынки предсказаний}

Четвёртый источник информации (дополнительный, на Stage 1.5) --- агент \texttt{ForesightCollector}, собирающий данные из рынков предсказаний для обогащения контекста прогноза.

% Основано на инженерной реализации и документации

\subsection{Polymarket CLOB API}

Основной источник данных --- \emph{Polymarket} (децентрализованная платформа предсказывающих рынков, построенная на блокчейне). Доступ осуществляется через HTTP REST API без требования аутентификации.

\textbf{Ключевые параметры API}:
\begin{itemize}
    \item URL: \texttt{https://clob.polymarket.com}
    \item Endpoints: \texttt{/markets}, \texttt{/prices}, \texttt{/trades}
    \item Rate limit: ~200 запросов/10 секунд
    \item Timeout: 10 секунд
\end{itemize}

Присоединение к рынкам выполняется через \texttt{conditionId} --- уникальный идентификатор условия события на блокчейне. Например, рынок «Will Trump win 2024 election?» может иметь conditionId = \texttt{0x1234...}. Этот ID используется как primary join key при агрегировании данных.

% Основано на Phase 6 completed

\subsection{Сбираемые метрики рынка}

Для каждого релевантного рынка собираются следующие метрики:

\begin{table}[h!]
\centering
\caption{Метрики Polymarket CLOB API}
\label{tab:polymarket_metrics}
\small
\begin{tabular}{|l|l|l|}
\hline
\textbf{Метрика} & \textbf{Источник} & \textbf{Назначение} \\
\hline
\texttt{conditionId} & markets endpoint & Уникальный идентификатор \\
\texttt{question} & markets & Формулировка предсказания \\
\texttt{volume} & prices & Объём торговли (всего) \\
\texttt{open\_interest} & prices & Открытый интерес (активные позиции) \\
\texttt{bid\_price} & prices & Лучшая цена покупки (YES) \\
\texttt{ask\_price} & prices & Лучшая цена продажи (YES) \\
\texttt{mid\_price} & computed & (bid + ask) / 2 \\
\texttt{spread} & computed & ask - bid (индикатор ликвидности) \\
\texttt{resolution\_date} & markets & Дата закрытия рынка \\
\texttt{last\_updated} & prices & Timestamp последнего обновления \\
\hline
\end{tabular}
\end{table}

\textit{Spread} (разница между ценами покупки и продажи) служит индикатором ликвидности рынка: узкий spread указывает на активную торговлю и уверенность участников.

% Основано на инженерной реализации

\subsection{Graceful degradation}

\textbf{Metaculus API} (альтернативный источник) в настоящее время недоступен из-за требования BENCHMARKING-уровня доступа (статус: запрошено 2026-03-29). Код пайплайна готов к интеграции, но на практике возвращает пустой список при попытке доступа.

Если Polymarket API недоступен (HTTP 5xx, timeout, rate limit):
\begin{enumerate}
    \item Retry с exponential backoff (3 попытки, макс. 30 сек)
    \item При повторных ошибках: логирование warning и возврат пустого списка `[]`
    \item Pipeline продолжает работу без данных рынков (деградация качества, но не критично)
\end{enumerate}

Успешный API-вызов возвращает список рынков с метриками. Релевантность рынка определяется совпадением между question/description рынка и событиями из EventCalendar (косинусное сходство эмбеддингов или простое ключевое слово matching).

---

\section{Идентификация событий: HDBSCAN-кластеризация}

Stage 2 пайплайна агломерирует 100--200 сигналов и 3--10 запланированных событий в 20 структурированных событийных нитей через плотностную кластеризацию.

% Source: 04-analysts.md:785-836

\subsection{Логика процесса}

Процесс разбивается на следующие этапы:

\begin{enumerate}
    \item \textbf{Подготовка текстов}: конкатенация title + summary для каждого SignalRecord

    \item \textbf{Эмбеддинг}: batch-эмбеддинг через LLM API (OpenAI embeddings или Voyage AI). Результат: матрица размером (N, 1536), где N = количество сигналов

    \item \textbf{Кластеризация HDBSCAN}: плотностная кластеризация эмбеддингов в пространстве косинусного сходства

    \item \textbf{LLM-лейблинг}: для каждого кластера генерируется название и summary через LLM

    \item \textbf{Интеграция ScheduledEvent}: запланированные события привязываются к кластерам на основе семантического сходства

    \item \textbf{Scoring}: расчёт \texttt{significance\_score} по многофакторной формуле

    \item \textbf{Ранжирование}: отбор top-20 по significance\_score

    \item \textbf{Trajectory Analysis}: для каждой нити анализируется текущее состояние, моментум, три сценария развития

    \item \textbf{Cross-Impact Matrix}: оценка перекрёстных влияний между событийными нитями
\end{enumerate}

% Source: 04-analysts.md:804-818

\subsection{Гиперпараметры HDBSCAN}

\begin{table}[h!]
\centering
\caption{Гиперпараметры кластеризации}
\label{tab:hdbscan_params}
\begin{tabular}{|l|l|l|}
\hline
\textbf{Параметр} & \textbf{Значение} & \textbf{Интерпретация} \\
\hline
\texttt{min\_cluster\_size} & 3 & Минимальный размер кластера (события, упомянутые в 3+ сигналах) \\
\texttt{min\_samples} & 2 & Минимальный размер neighbourhood для плотности \\
\texttt{metric} & cosine & Метрика расстояния в пространстве эмбеддингов \\
\hline
\end{tabular}
\end{table}

\texttt{min\_cluster\_size} = 3 означает, что кластер формируется только если он содержит минимум 3 сигнала с близкими эмбеддингами (косинусное сходство). Это фильтрует шум (случайные совпадения одного-двух заголовков) и выделяет только события, получившие достаточное освещение в множественных источниках.

% Source: 04-analysts.md:812-818

\subsection{Формула значимости события}

Значимость события рассчитывается как взвешенная сумма пяти компонент:

\begin{equation}
\text{significance\_score} = 0.30 \times \text{importance} + 0.25 \times \text{cluster\_size\_norm} + 0.20 \times \text{recency} + 0.15 \times \text{source\_diversity} + 0.10 \times \text{entity\_prominence}
\label{eq:significance}
\end{equation}

Компоненты определяются следующим образом:

\begin{enumerate}
    \item \textbf{importance} (вес 0.30) --- оценка LLM медийной важности события в контексте выбранного издания (0.0--1.0). Высокое значение для событий, которые гарантированно будут в топе новостей.

    \item \textbf{cluster\_size\_norm} (вес 0.25) --- нормализованный размер кластера. Если кластер содержит $k$ сигналов, а максимальный размер $k_{\max}$, то этот компонент = $k / k_{\max}$ (0.0--1.0). Большие кластеры указывают на высокую новостную активность.

    \item \textbf{recency} (вес 0.20) --- свежесть сигналов в кластере. Сигналы, опубликованные в последние 24 часа, получают полный балл 1.0. Более старые сигналы штрафуются экспоненциально: $\exp(-t / \tau)$, где $t$ --- возраст (дни), $\tau = 3$ (временная постоянная).

    \item \textbf{source\_diversity} (вес 0.15) --- количество уникальных источников в кластере, нормализованное: $\min(1.0, |\text{sources}| / 5)$. Событие, покрытое 5+ различными источниками, получает максимальный балл.

    \item \textbf{entity\_prominence} (вес 0.10) --- число упоминаний высокопрофильных сущностей (главы государств, ЕЦБ, ООН и т.д.) в сигналах кластера, нормализованное к [0, 1]. Для типичного кластера этот компонент невелик, но может быть решающим для крупных политических событий.
\end{enumerate}

% Source: 04-analysts.md:769-836

\subsection{Выход: EventThread[]}

Результат кластеризации --- список объектов \texttt{EventThread} (до 20 записей), упорядоченный по \texttt{significance\_score}:

\begin{table}[h!]
\centering
\caption{Поля EventThread}
\label{tab:event_thread}
\footnotesize
\begin{tabular}{|l|l|l|}
\hline
\textbf{Поле} & \textbf{Тип} & \textbf{Содержание} \\
\hline
\texttt{id} & str & Уникальный идентификатор нити \\
\texttt{title} & str & Название события (сгенерировано LLM из кластера) \\
\texttt{summary} & str & Summary событийной нити \\
\texttt{signal\_ids} & list[str] & ID входящих сигналов \\
\texttt{scheduled\_events} & list[str] & ID привязанных запланированных событий \\
\texttt{cluster\_size} & int & Количество сигналов в кластере \\
\texttt{significance\_score} & float & Итоговая значимость (0.0--1.0) \\
\texttt{sources} & list[str] & Уникальные источники сигналов \\
\texttt{entities} & list[str] & Агрегированные именованные сущности \\
\texttt{published\_at} & datetime & Время публикации самого старого сигнала \\
\texttt{last\_updated} & datetime & Время самого свежего сигнала \\
\texttt{trajectory} & EventTrajectory | None & Анализ траектории (Stage 3) \\
\hline
\end{tabular}
\end{table}

Каждая EventThread представляет связанный набор новостей об одном событии или процессе, готовый для анализа на последующих стадиях пайплайна (Stage 3: тренд-анализ, Stage 4--5: Дельфи-прогнозирование).

---

\subsection{Обработка ошибок и граничные случаи}

\begin{table}[h!]
\centering
\caption{Стратегии обработки ошибок в Stage 2}
\label{tab:stage2_errors}
\footnotesize
\begin{tabular}{|l|l|}
\hline
\textbf{Ситуация} & \textbf{Действие} \\
\hline
< 10 сигналов на входе & Вернуть как есть + warning; HDBSCAN может не сформировать кластеры \\
Эмбеддинг API unavailable & Fallback: использовать простую текстовую дистанцию (Jaccard) \\
0 кластеров найдено & Трактовать каждый сигнал как отдельный EventThread \\
LLM не парсит JSON при лейблинге & Retry 1 раз; использовать generic title при неудаче \\
\hline
\end{tabular}
\end{table}



\newpage
% =====================================================================
% PARTS III + IV: Analysis + Delphi Method (Stages 3-5)
% =====================================================================
% =====================================================================
% ЧАСТЬ III: Анализ траекторий и сценариев
% =====================================================================

\section{Три параллельных аналитика (Стадия 3)}

% Source: docs/04-analysts.md, Stage 3 specification

На Стадии 3 события, выделенные EventTrendAnalyzer, поступают в три специализированные аналитические агента, работающие \emph{параллельно}:

\begin{enumerate}
    \item \textbf{Геополитический аналитик} (\emph{Geopolitical Analyst}) — анализ стратегических акторов, баланса сил, альянсов, вероятности эскалации;
    \item \textbf{Экономический аналитик} (\emph{Economic Analyst}) — анализ финансовых потоков, рыночных индикаторов, цепочек поставок, сценариев воздействия на индексы;
    \item \textbf{Медиа-аналитик} (\emph{Media Analyst}) — оценка медийной ценности события через шести-мерный фреймворк \emph{Galtung \& Ruge} (1965).
\end{enumerate}

Минимум успешных оценок: $\text{min\_successful} = 2$ из 3. Архитектура допускает отказ одного аналитика без отмены Дельфи-прогноза.

\subsection{Геополитический анализ}

Для каждого события строится профиль ключевых стратегических акторов:
\begin{itemize}
    \item \textbf{Стратегические акторы}: государства, альянсы (НАТО, ШОС), международные организации (ООН, ЕС).
    \item \textbf{Баланс сил}: относительная мощь, направления сдвига в расстановке (кто усиливается, кто ослабевает).
    \item \textbf{Вероятность эскалации}: для конфликтов, дипломатических кризисов — числовая оценка (0–1).
    \item \textbf{Эффекты второго порядка}: цепь причин-следствий, как действие одного актора спровоцирует ответ другого.
\end{itemize}

Аналитическая рамка сочетает неореализм (структура системы определяет поведение акторов) с конструктивизмом (восприятие угроз зависит от идентичности и нарратива).

\subsection{Экономический анализ}

Для каждого события:
\begin{itemize}
    \item \textbf{Затронутые индикаторы}: валютные пары, цены на сырьё, фондовые индексы, облигационные спреды.
    \item \textbf{Фискальные последствия}: взаимосвязь с бюджетной политикой государств.
    \item \textbf{Цепочки поставок}: географические разрывы, сектора, уязвимые к эскалации или стабилизации.
    \item \textbf{Рыночный сигнал}: что рынки уже закладывают в цены и что ещё является сюрпризом.
\end{itemize}

Ключевой принцип: \emph{«Follow the money»} — экономические стимулы раскрывают истинные намерения и ограничения акторов.

\subsection{Медиа-анализ: шести-мерный фреймворк новостной ценности}

% Source: Galtung & Ruge 1965, Entman 1993 framing theory

Новостная ценность события оценивается по шести критериям (Galtung \& Ruge, 1965):

\begin{enumerate}
    \item \textbf{Timeliness} (своевременность): событие происходит сейчас, не архивное.
    \item \textbf{Magnitude} (масштаб): велико ли событие по количеству затронутых лиц или активов.
    \item \textbf{Prominence} (известность): вовлечены ли топ-фигуры, знакомые аудитории.
    \item \textbf{Proximity} (близость): географическая, культурная, тематическая близость к аудитории издания.
    \item \textbf{Conflict} (конфликтность): наличие явного противостояния, эмоционального заряда.
    \item \textbf{Novelty} (новизна): событие неожиданно, радикально отличается от прецедентов.
\end{enumerate}

Кроме того, медиа-аналитик оценивает:
\begin{itemize}
    \item \textbf{Редакционный fit} (0–1): соответствие события редакционной линии издания.
    \item \textbf{Медиа-насыщенность} (0–1): как долго тема находилась в топе новостей. Если > 14 дней подряд — редакция ищет свежий угол или отстранится.
    \item \textbf{Позиция в цикле}: находится ли тема на пике внимания, идёт на спад или восходящий тренд.
\end{itemize}

Выход: \texttt{MediaAssessment} с итоговым баллом новостной ценности и прогнозом вероятности попадания в выпуск данного издания.

---

% =====================================================================
% МОДЕЛИРОВАНИЕ ТРАЕКТОРИЙ СОБЫТИЙ
% =====================================================================

\section{Моделирование траекторий событий (Стадия 3)}

% Source: docs/04-analysts.md, EventTrajectory schema

После параллельного анализа для каждого события строится модель его траектории развития. Траектория описывает:

\begin{enumerate}
    \item \textbf{Текущее состояние} — где находится событие сейчас (описание 2–3 предложениями).
    \item \textbf{Моментум} (\emph{momentum}) — вектор развития события.
    \item \textbf{Три сценария} — BASELINE, OPTIMISTIC/PESSIMISTIC, WILDCARD.
    \item \textbf{Ключевые факторы} — 3–5 движущих сил, определяющих развитие.
    \item \textbf{Неопределённости} — 2–3 основные точки неопределённости.
\end{enumerate}

\subsection{Моментум события}

Моментум описывает вектор развития событийной нити в краткосрочной перспективе:

\begin{itemize}
    \item \textbf{Escalating} — ситуация накаляется, вероятность критических событий растёт.
    \item \textbf{Stable} — статус-кво поддерживается, существенных изменений не ожидается.
    \item \textbf{De-escalating} — напряжённость снижается, кризис миновал пик.
    \item \textbf{Emerging} — новое событие только начинает появляться в медийном пространстве.
    \item \textbf{Culminating} — событие приближается к критической точке, решению или кульминации.
    \item \textbf{Fading} — событие теряет актуальность, медийное внимание падает.
\end{itemize}

\subsection{Три сценария развития}

Для каждого события определяются три сценария с указанием:
\begin{itemize}
    \item Краткое описание (2–3 предложения).
    \item Назначенная вероятность (0.0–1.0), сумма по трём сценариям = 1.0.
    \item 2–3 ключевых индикатора, указывающих на реализацию этого сценария.
    \item Потенциальный заголовок, который может породить сценарий.
\end{itemize}

\textbf{Типы сценариев}:
\begin{enumerate}
    \item \textbf{BASELINE} — наиболее вероятное развитие событий на основе текущего моментума.
    \item \textbf{OPTIMISTIC} или \textbf{PESSIMISTIC} — улучшение или ухудшение ситуации по ключевым параметрам.
    \item \textbf{WILDCARD} — неожиданный поворот, требующий активации какого-либо из выявленных рисков.
\end{enumerate}

\subsection{Матрица перекрёстных влияний}

% Source: docs/04-analysts.md, CrossImpactMatrix schema

События в информационном пространстве не независимы. Развитие одного события влияет на вероятность других.

Матрица перекрёстных влияний строится как разреженное представление:
\begin{equation}
    \text{CrossImpactEntry:} \quad (source\_id, target\_id, impact\_score) \in [-1, 1] \times \mathbb{R}
\end{equation}

где:
\begin{itemize}
    \item $source\_id$ — событие-причина.
    \item $target\_id$ — событие-следствие.
    \item $impact\_score \in [-1, 1]$ — сила и направление влияния:
    \begin{itemize}
        \item $> 0$ — усиливающее влияние (источник ускоряет, вероятнее делает целевое событие).
        \item $< 0$ — ослабляющее влияние (источник замедляет, маловероятнее делает целевое).
        \item $= 0$ — отсутствие влияния (не включается в матрицу).
    \end{itemize}
\end{itemize}

Для типичного портфеля из 20 событий полная матрица содержала бы $20 \times 19 = 380$ пар. На практике значимых связей 30–50, остальные игнорируются. Эта разреженность критична для масштабируемости Дельфи при большом числе событий.

---

% =====================================================================
% ЧАСТЬ IV: МЕТОД ДЕЛЬФИ ДЛЯ LLM-АГЕНТОВ
% =====================================================================

\section{Пять экспертных персон (Дельфи)}

% Source: src/agents/forecasters/personas.py, lines 24-242

Дельфи-система построена на пяти независимых \emph{экспертных персонах} (\emph{expert personas}) — LLM-агентов с чётко определёнными аналитическими рамками, когнитивными смещениями и начальными весами. Каждая персона работает на одной и той же модели (\texttt{claude-opus-4.6}), но различается \emph{системным промптом}, определяющим её идентичность, аналитическую методологию и осознанные когнитивные смещения.

\begin{table}[h]
\small
\centering
\begin{tabular}{|c|l|l|c|l|}
\hline
\textbf{ID} & \textbf{Имя} & \textbf{Роль} & \textbf{Вес} & \textbf{Основная методология} \\
\hline
\hline
1 & REALIST & Аналитик рисков & 0.22 & Базовые ставки, Tetlock \emph{Superforecasting} \\
\hline
2 & GEOSTRATEG & Специалист МО & 0.20 & Неореализм, cui bono, баланс сил \\
\hline
3 & ECONOMIST & Макроэкономист & 0.20 & Follow the money, рыночные индикаторы \\
\hline
4 & MEDIA\_EXPERT & Редактор/медиа-аналитик & 0.18 & Гейткипинг \emph{White 1950}, фрейминг \emph{Entman 1993} \\
\hline
5 & DEVILS\_ADVOCATE & Контрариан/red teamer & 0.20 & Pre-mortem, \emph{Black swan} detection, Taleb \\
\hline
\end{tabular}
\caption{Экспертные персоны Дельфи: идентичность, методология и начальные веса}
\label{table:personas}
\end{table}

Суммарный вес: $0.22 + 0.20 + 0.20 + 0.18 + 0.20 = 1.00$.

\subsection{Персона 1: Реалист-аналитик}

\textbf{Идентичность}: Опытный аналитик политических рисков (условный профиль: Eurasia Group, Oxford Analytica) с 20-летним опытом работы в консалтинге. Мыслит категориями базовых ставок, исторических прецедентов и институциональной инерции.

\textbf{Аналитическая рамка}:
\begin{enumerate}
    \item \emph{Базовые ставки} — первый вопрос: «Как часто подобное происходило за 10–20 лет? Каковы были исходы?»
    \item \emph{Институциональная инерция} — системы меняются медленнее, чем кажется.
    \item \emph{Outside view перед inside view} — начинать с внешнего вида (Tetlock, Superforecasting).
    \item \emph{Калибровка над убеждённостью} — лучше неуверенная точная оценка, чем уверенная ошибочная.
    \item \emph{Разделение вопросов} — «Произойдёт ли?» и «Попадёт ли в заголовки?» — разные задачи.
\end{enumerate}

\textbf{Когнитивное смещение} (управляемое):
\begin{itemize}
    \item \textbf{Переоценивает}: инерция статус-кво, предсказуемость бюрократических процессов.
    \item \textbf{Недооценивает}: чёрные лебеди, скорость эскалации, роль личных решений лидеров.
    \item \textbf{Якорь при неопределённости}: исторические прецеденты.
\end{itemize}

\textbf{Начальный вес}: 0.22 (повышенный — базовые ставки исторически хорошо калиброваны).

\subsection{Персона 2: Геополитический стратег}

\textbf{Идентичность}: Специалист по международным отношениям (IISS, Chatham House, ИМЭМО РАН) с фокусом на силовые балансы, альянсы и стратегические интересы государств.

\textbf{Аналитическая рамка}:
\begin{enumerate}
    \item \emph{Cui bono} — кто выигрывает, кто проигрывает?
    \item \emph{Деревья решений} — эффекты второго и третьего порядка.
    \item \emph{Неореализм + конструктивизм} — структура системы определяет поведение; идентичности и нарративы формируют восприятие.
    \item \emph{Красные линии и пороги} — различение реальных красных линий от блефа.
\end{enumerate}

\textbf{Когнитивное смещение} (управляемое):
\begin{itemize}
    \item \textbf{Переоценивает}: рациональность государственных акторов, значимость геополитических факторов.
    \item \textbf{Недооценивает}: внутриполитические случайности, роль технологий, экономическую мотивацию негосударственных акторов.
    \item \textbf{Якорь}: модель великодержавного соперничества.
\end{itemize}

\textbf{Начальный вес}: 0.20.

\subsection{Персона 3: Экономический аналитик}

\textbf{Идентичность}: Макроэкономист и рыночный аналитик (Goldman Sachs Research, The Economist Intelligence Unit) с фокусом на потоки капитала, фискальную политику, цены на сырьё и корпоративные интересы.

\textbf{Аналитическая рамка}:
\begin{enumerate}
    \item \emph{Follow the money} — базовый принцип.
    \item \emph{Рациональный актор с бюджетными ограничениями} — экономические стимулы раскрывают истинные намерения.
    \item \emph{Экономический календарь} — источник предсказуемых новостных поводов (релизы, заседания ЦБ).
    \item \emph{Рыночные сигналы} содержат информацию (кривая доходностей, CDS спреды, цена золота).
\end{enumerate}

\textbf{Когнитивное смещение} (управляемое):
\begin{itemize}
    \item \textbf{Переоценивает}: экономическую рациональность, предсказательную силу рыночных индикаторов.
    \item \textbf{Недооценивает}: идеологические факторы, культурные войны, события, обусловленные эмоциями.
    \item \textbf{Якорь}: экономический календарь и рыночный консенсус.
\end{itemize}

\textbf{Начальный вес}: 0.20.

\subsection{Персона 4: Медиа-эксперт}

\textbf{Идентичность}: Бывший редактор крупного новостного агентства (Reuters, ТАСС) + преподаватель журналистики. Понимает внутреннюю кухню редакции: дедлайны, источники, редакционную политику, конкуренцию за внимание.

\textbf{Аналитическая рамка}:
\begin{enumerate}
    \item \emph{Гейткипинг} (\textbf{White}, 1950) — не всё важное публикуется. Редакционные фильтры ограничивают пропускную способность.
    \item \emph{Фрейминг} (\textbf{Entman}, 1993) — одно событие, десять подач. Выбор фрейма определяет заголовок и угол.
    \item \emph{Медиа-насыщенность} — тема в топе 14+ дней подряд $\to$ издание ищет свежий угол или отстранится.
    \item \emph{Конкуренция за внимание} — параллельные события конкурируют за аудиторию.
\end{enumerate}

\textbf{Когнитивное смещение} (управляемое):
\begin{itemize}
    \item \textbf{Переоценивает}: значимость медиа-циклов и информационных трендов, предсказуемость редакционных решений.
    \item \textbf{Недооценивает}: технические и процедурные события, медленно развивающиеся кризисы.
    \item \textbf{Якорь}: текущий новостной цикл и тематический баланс издания.
\end{itemize}

\textbf{Начальный вес}: 0.18 (ниже среднего — медиа-экспертиза важна для формулировки заголовков, менее для вероятностных оценок).

\subsection{Персона 5: Адвокат дьявола}

\textbf{Идентичность}: Систематический контрариан и специалист по анализу рисков (red team, разведывательное сообщество, школа Nassim Taleb). Задача — найти то, что остальные пропустили.

\textbf{Аналитическая рамка}:
\begin{enumerate}
    \item \emph{Pre-mortem} — «Представь, прогноз провалился. Что именно пошло не так?» (Klein, 1989).
    \item \emph{Steelmanning} — строить максимально сильную версию противоположной позиции, потом атаковать.
    \item \emph{Каскадные зависимости} — если A спорно и B зависит от A, то цепочка становится хрупкой.
    \item \emph{Чёрные лебеди} (Талеб) — маловероятные, но высокоимпактные сценарии.
\end{enumerate}

\textbf{Когнитивное смещение} (управляемое, и в этом случае — желаемое):
\begin{itemize}
    \item \textbf{Переоценивает}: вероятность чёрных лебедей, хрупкость систем, неочевидные причинно-следственные связи.
    \item \textbf{Недооценивает}: устойчивость статус-кво, вероятность «скучных» исходов, стабильность институтов.
    \item \textbf{Якорь}: маловероятные, но высокоимпактные сценарии.
\end{itemize}

\textbf{Начальный вес}: 0.20 (парадоксально высокий — контрарианские прогнозы добавляют ценность ансамблю при агрегации, даже если редко верны в большинстве случаев).

---

% =====================================================================
% РАУНД 1: НЕЗАВИСИМЫЕ ОЦЕНКИ
% =====================================================================

\section{Раунд 1: независимые оценки пяти персон}

% Source: src/agents/forecasters/personas.py, lines 269-333, and docs/05-delphi-pipeline.md §2

Все пять персон запускаются \emph{параллельно}. Каждая получает:

\begin{itemize}
    \item \textbf{Профиль издания} (\texttt{OutletProfile}): географический охват, целевая аудитория, редакционная политика, тематический баланс.
    \item \textbf{Траектории событий} (\texttt{EventTrajectory[]}): до 20 событий с текущим состоянием, моментумом, сценариями, ключевыми факторами.
    \item \textbf{Матрица перекрёстных влияний} (\texttt{CrossImpactMatrix}): разреженное представление зависимостей между событиями.
    \item \textbf{Система промпта} — специфичный для персоны системный промпт, определяющий её идентичность, методологию и смещения.
\end{itemize}

Параметры LLM для R1:
\begin{itemize}
    \item \textbf{Модель}: \texttt{anthropic/claude-opus-4.6}
    \item \textbf{Temperature}: $T = 0.7$ (некоторая стохастичность, но управляемая)
    \item \textbf{max\_tokens}: 4096
    \item \textbf{json\_mode}: истина (структурированный выход)
\end{itemize}

Минимум успешных агентов: $\text{min\_successful} = 3$ из 5. При отказе одного агента раунд продолжается.

\subsection{Выходная схема: PersonaAssessment}

Каждая персона возвращает структурированную оценку \texttt{PersonaAssessment}:

\begin{definition}[\texttt{PersonaAssessment} — оценка от одной персоны]
\begin{verbatim}
{
  persona_id: str,               # "realist", "geostrateg", и т.д.
  round_number: int,             # 1 или 2
  predictions: PredictionItem[], # 5–15 прогнозов
  cross_impacts_noted: str[],    # замеченные перекрёстные влияния
  blind_spots: str[],            # что группа может пропустить
  confidence_self_assessment: float,  # (0–1) мета-уверенность

  # Только для раунда 2:
  revisions_made: str[],         # что было изменено
  revision_rationale: str        # почему
}
\end{verbatim}
\end{definition}

\subsection{Выходная схема: PredictionItem}

Каждый элемент \texttt{predictions[]} содержит одно предсказание события:

\begin{definition}[\texttt{PredictionItem} — единичный прогноз]
\begin{verbatim}
{
  event_thread_id: str,              # ID события из Stage 3
  prediction: str,                   # конкретное утверждение
  probability: float,                # вероятность (0.0–1.0, не кратна 5%)
  newsworthiness: float,             # (0–1) вероятность попадания в выпуск
  scenario_type: ScenarioType,       # BASELINE | OPTIMISTIC | PESSIMISTIC | WILDCARD
  reasoning: str,                    # цепочка рассуждений (3–7 предложений)
  key_assumptions: str[],            # 2–4 ключевые предпосылки
  evidence: str[],                   # ссылки на факты из входных данных
  conditional_on: str[]              # ID других прогнозов, от которых зависит этот
}
\end{verbatim}
\end{definition}

\textbf{Критические требования}:
\begin{enumerate}
    \item \textbf{Точность вероятностей}: каждая вероятность должна быть уникальна, не кратна 5 или 10 (0.63, не 0.60). Это требование (Tetlock) улучшает калибровку.
    \item \textbf{Диапазон}: $0.03 \leq \text{probability} \leq 0.97$. Вероятности 0.00 или 1.00 запрещены (они означают избыточную уверенность).
    \item \textbf{Обоснование}: каждый прогноз должен иметь явную цепочку рассуждений и ключевые предпосылки, на которых основана оценка.
\end{enumerate}

---

% =====================================================================
% МЕДИАТОР: СИНТЕЗ РАЗНОГЛАСИЙ
% =====================================================================

\section{Медиатор: синтез разногласий (Стадия 5a)}

% Source: src/agents/forecasters/mediator.py, lines 1-269; docs/prompts/mediator.md §1-10

После Раунда 1 пять независимых оценок поступают в \emph{Медиатор} — специализированный агент, чья задача \emph{не прогнозировать}, а \emph{структурировать разногласия} для содержательного второго раунда.

\subsection{Классический Дельфи vs. LLM-Дельфи}

В классическом методе Дельфи (Dalkey \& Helmer, 1963) обратная связь для второго раунда состоит из агрегированной статистики: медиана, квартили, гистограмма.

Однако исследование \textbf{DeLLMphi} (Zhao et al., 2024) показало: если LLM-агентам дать только медианные оценки группы, они слегка сдвигаются к ней без содержательной ревизии аргументов. Это явление называется \emph{Degeneration-of-Thought} (Liang et al., 2024, EMNLP).

\textbf{Наше решение}: вместо голой статистики, Медиатор формулирует \emph{содержательные конкретные вопросы}, которые раскрывают фактические разногласия и заставляют агентов пересматривать не числа, а аргументы.

\subsection{Три этапа медиации}

\subsubsection{Этап 1: Алгоритмическая классификация событий}

Для каждого события Медиатор собирает все оценки пяти персон и вычисляет:

\begin{enumerate}
    \item \textbf{Разброс} (\emph{spread}): $\text{spread} = \max(\text{probabilities}) - \min(\text{probabilities})$
    \item \textbf{Медиана}: $\text{median}(\text{probabilities})$
    \item \textbf{Число агентов, упомянувших событие}: $n \in [0, 5]$
\end{enumerate}

События классифицируются как:

\begin{enumerate}
    \item \textbf{Консенсус} (\emph{Consensus Area}): $\text{spread} < 0.15$ \textbf{AND} $n \geq 3$

    \textit{Интерпретация}: эксперты согласны, нет необходимости в ревизии. Это событие попадёт в финальный прогноз с доверием.

    \item \textbf{Разногласие} (\emph{Dispute Area}): $\text{spread} \geq 0.15$

    \textit{Интерпретация}: эксперты не согласны. Медиатор формулирует конкретный фактический вопрос, ответ на который мог бы сблизить позиции.

    \item \textbf{Пробел} (\emph{Gap Area}): $n < 3$

    \textit{Интерпретация}: события упомянуты менее чем тремя экспертами. Потенциально важные события, которые группа в целом упустила.
\end{enumerate}

\subsubsection{Этап 2: Проверка каскадных зависимостей}

\emph{CrossImpactFlag}: если прогноз события A зависит от события B (через поле \texttt{conditional\_on}), и по B существует разногласие, то это создаёт цепную неопределённость, которую нужно подсветить.

\subsubsection{Этап 3: LLM-обогащение синтеза}

Медиатор получает:
\begin{itemize}
    \item Анонимизированные оценки (метки Эксперт A–E, см. ниже).
    \item Траектории событий (для контекста).
    \item Результаты этапов 1–2 (алгоритмическая классификация).
\end{itemize}

И через LLM-вызов \emph{обогащает} синтез:
\begin{itemize}
    \item Для каждого разногласия: формулирует one-liner ключевого вопроса, ответ на который сближает позиции.
    \item Для каждого пробела: объясняет, почему это может быть важным упущением.
    \item Генерирует общее резюме (2–3 предложения): сколько консенсусов, расхождений, пробелов.
\end{itemize}

Параметры LLM для Медиатора:
\begin{itemize}
    \item \textbf{Модель}: \texttt{anthropic/claude-opus-4.6}
    \item \textbf{Temperature}: 0.7 (нейтральная)
    \item \textbf{json\_mode}: истина
\end{itemize}

\subsection{Анонимность как защита от группового мышления}

% Source: src/agents/forecasters/mediator.py, lines 181-206; docs/prompts/mediator.md §7-9

Критический механизм защиты от давления большинства (Zhang et al., 2024, ACL): **метки экспертов переиграются случайно при каждом запуске**.

Алгоритм:
\begin{enumerate}
    \item Для каждого набора оценок генерируется детерминированное дерево случайных чисел, основанное на хеше идентификаторов персон.
    \item Метки $\{$Эксперт A, B, C, D, E$\}$ перемешиваются.
    \item Одна и та же персона получает разную метку при каждом прогоне.
\end{enumerate}

Эффект: агент не может узнать себя по метке и подстраиться под неё. Это предотвращает конформность и сохраняет независимость позиций меньшинства.

\subsection{Выходная схема: MediatorSynthesis}

\begin{definition}[\texttt{MediatorSynthesis} — структурированный синтез разногласий]
\begin{verbatim}
{
  consensus_areas: ConsensusArea[],     # spread < 0.15, n >= 3
  disputes: DisputeArea[],              # spread >= 0.15 + ключевой вопрос
  gaps: GapArea[],                      # n < 3
  cross_impact_flags: CrossImpactFlag[], # зависимости между спорными событиями
  overall_summary: str,                 # 2–3 предложения об итоговой картине
  supplementary_facts: str[]            # заполняется позже, если нужен supervisor_search
}
\end{verbatim}
\end{definition}

где:

\begin{definition}[\texttt{ConsensusArea}]
\begin{verbatim}
{
  event_thread_id: str,
  median_probability: float,  # (0–1)
  spread: float,              # < 0.15
  num_agents: int             # >= 3
}
\end{verbatim}
\end{definition}

\begin{definition}[\texttt{DisputeArea}]
\begin{verbatim}
{
  event_thread_id: str,
  median_probability: float,
  spread: float,              # >= 0.15
  positions: AnonymizedPosition[],  # позиции каждого эксперта
  key_question: str           # конкретный фактический вопрос для R2
}

AnonymizedPosition: {
  agent_label: str,           # "Эксперт A", "Эксперт B", и т.д.
  probability: float,
  reasoning_summary: str,     # первые 200 символов обоснования
  key_assumptions: str[]
}
\end{verbatim}
\end{definition}

\begin{definition}[\texttt{GapArea}]
\begin{verbatim}
{
  event_thread_id: str,
  mentioned_by: str[],        # список меток экспертов, упомянувших событие
  note: str                   # почему это может быть важным пробелом
}
\end{verbatim}
\end{definition}

\begin{definition}[\texttt{CrossImpactFlag}]
\begin{verbatim}
{
  prediction_event_id: str,    # событие A
  depends_on_event_id: str,    # событие B (спорное)
  note: str                    # описание зависимости A$\to$B
}
\end{verbatim}
\end{definition}

---

% =====================================================================
% РАУНД 2: РЕВИЗИЯ С ОБРАТНОЙ СВЯЗЬЮ
% =====================================================================

\section{Раунд 2: ревизия с обратной связью}

% Source: src/agents/forecasters/personas.py, lines 269-308; docs/05-delphi-pipeline.md §3.2-3.3

Все пять персон запускаются \emph{снова параллельно}. Каждая получает:

\begin{itemize}
    \item \textbf{Свои собственные оценки из R1}.
    \item \textbf{MediatorSynthesis} — анонимизированные позиции других экспертов и ключевые вопросы для разногласий.
    \item \textbf{Гвард независимости} (\emph{Independence Guard}): явная инструкция — \textit{«НЕ сдвигай числа просто потому, что все думают иначе. Ответь на ключевой вопрос медиатора аргументированно»}.
\end{itemize}

Параметры LLM для R2:
\begin{itemize}
    \item \textbf{Модель}: \texttt{anthropic/claude-opus-4.6}
    \item \textbf{Temperature}: $T = 0.6$ (снижена с 0.7, чтобы уменьшить шум, но сохранить творчество)
    \item \textbf{max\_tokens}: 4096
    \item \textbf{json\_mode}: истина
\end{itemize}

Минимум успешных агентов: $\text{min\_successful} = 3$ из 5.

\subsection{Выходная схема: RevisedAssessment}

Структура идентична \texttt{PersonaAssessment}, но с дополнительными полями:

\begin{verbatim}
{
  # Все поля как в R1...
  predictions: PredictionItem[],
  confidence_self_assessment: float,

  # Плюс, для R2:
  revisions_made: str[],        # список ревизий
  revision_rationale: str       # почему эти ревизии
}
\end{verbatim}

Примеры ревизий:
\begin{itemize}
    \item «Повысил вероятность с 0.42 до 0.58 — ключевой вопрос медиатора о венгерском вето важен, я пересмотрел рассуждения».
    \item «Сохранил 0.25 — позиция меньшинства, но фактический вопрос не рассеял мою неопределённость».
    \item «Добавил новый прогноз на событие, упомянутое одним экспертом — теперь вижу важность этого пробела».
\end{itemize}

---

% =====================================================================
% ТЕОРЕТИЧЕСКОЕ ОБОСНОВАНИЕ МЕТОДА ДЕЛЬФИ
% =====================================================================

\section{Теоретическое обоснование метода Дельфи для LLM-агентов}

\subsection{Классический Дельфи и его преимущества}

\emph{Метод Дельфи} (Dalkey \& Helmer, RAND Corporation, 1963) — структурированная техника группового прогнозирования, основанная на четырёх принципах:

\begin{enumerate}
    \item \textbf{Анонимность}: эксперты не знают, кто дал какую оценку.
    \item \textbf{Итерация}: несколько раундов оценок с обратной связью.
    \item \textbf{Контролируемая обратная связь}: участники получают агрегированную статистику группы (медиана, квартили, гистограмма).
    \item \textbf{Статистическая агрегация}: финальный результат — усреднённая оценка группы.
\end{enumerate}

Доказанные преимущества (Rowe \& Wright, 2001, 2005):
\begin{itemize}
    \item Подавление эффекта якорения — эксперты сдвигают позиции на основе данных, а не авторитета.
    \item Снижение давления авторитета — анонимность защищает меньшинство.
    \item Более калиброванные вероятности — групповая оценка часто превосходит одиночных экспертов.
\end{itemize}

\subsection{Адаптация для LLM-агентов}

Наша реализация заменяет человеческих экспертов на LLM-агентов с чётко определёнными когнитивными профилями. Ключевые отличия:

\begin{enumerate}
    \item \textbf{Детерминизм} — каждый агент получает одинаковые входные данные (траектории, перекрёстные влияния) и системный промпт, определяющий его аналитическую рамку.

    \item \textbf{Отсутствие давления авторитета в классическом смысле} — агенты не знают о мнениях друг друга до медиации. Давление группы вводится контролируемо через медиатора, а не через выраженное большинство.

    \item \textbf{Воспроизводимость} — результаты должны быть воспроизводимы (с фиксированными семенами случайности).
\end{enumerate}

\subsection{Ключевые исследования}

\subsubsection{AIA Forecaster: необходимость модельного разнообразия}

Schoenegger et al. (2024) показали: если все агенты в ансамбле используют одну модель, их ошибки \emph{коррелируют}. Коррелированные ошибки не компенсируются при агрегации.

\textbf{Вывод}: ансамбль из 5 экземпляров одной модели лишь незначительно лучше одного вызова той же модели.

\textbf{Наше решение}: все пять персон работают на одной модели (\texttt{claude-opus-4.6}), но \emph{разнообразие ошибок обеспечивается система prompts, когнитивными смещениями и аналитическими рамками}, не различием моделей. Это позволяет:
\begin{itemize}
    \item Контролировать стоимость (все вызовы к одному провайдеру).
    \item Гарантировать минимальное качество рассуждений (Opus 4.6 — лучший в своём ценовом диапазоне).
    \item Достичь разнообразия через когнитивные профили (исследовано).
\end{itemize}

\subsubsection{DeLLMphi и проблема Degeneration-of-Thought}

Zhao et al. (2024) воспроизвели классический Дelphi с LLM и обнаружили критическую проблему:

\begin{enumerate}
    \item В R1 агенты дают независимые оценки.
    \item В R2 им показывают медианные оценки группы.
    \item Агенты \emph{слегка сдвигаются к медиане}, но \textbf{без содержательной ревизии аргументов}.
\end{enumerate}

Это явление названо \emph{Degeneration-of-Thought} (Liang et al., 2024, EMNLP): LLM, утвердившиеся в позиции, не способны к \emph{genuine revision} через self-reflection одной только статистики.

\textbf{Наше решение}: Медиатор не просто даёт медиану. Он формулирует \emph{конкретные фактические вопросы}, на которые необходимо ответить:

\begin{equation}
    \text{Ключевой вопрос:} \quad Q = \text{«Что именно нужно проверить для разрешения этого спора?»}
\end{equation}

Lorenz \& Fritz (2025, arXiv:2602.08889) показали: когда агентам задаются \emph{substantive questions}, корреляция их оценок с ground truth достигает $r = 0.87$--$0.95$. При простой статистике $r \approx 0.60$--$0.70$.

\subsubsection{Защита позиций меньшинства}

Zhang et al. (2024, ACL) показали: LLM-агенты воспроизводят human-like conformity с \emph{bandwagon score} 0.524 в GPT-3.5. То есть, когда большинство согласно, агент сдвигается к большинству даже если изначально думал иначе.

\textbf{Наше решение}: анонимность с ротацией меток. Метки (A–E) переиграются при каждом прогоне так, что агент не может узнать себя и подстраиться.

\subsubsection{Два раунда достаточно}

Классические исследования Дельфи (Rowe \& Wright, 2005) показали:
\begin{itemize}
    \item Раунд 1 $\to$ Раунд 2: существенное улучшение (до 80% сходимости).
    \item Раунд 2 $\to$ Раунд 3: маргинальное улучшение (дополнительно ~5–10%).
    \item Раунд 3+: убывающая отдача.
\end{itemize}

При LLM стоимость растёт линейно с числом раундов. \textbf{Наша архитектура: ровно 2 раунда} — баланс между качеством и стоимостью.

---

% =====================================================================
% Заключение логической архитектуры Дельфи
% =====================================================================

\section{Логическая архитектура двухраундового Дельфи}

Полный цикл:

\begin{equation}
\begin{aligned}
    \text{R1:} \quad &\text{5 personas} \xrightarrow{\parallel} \text{PersonaAssessment}[] \\
    \text{Mediation:} \quad &\text{PersonaAssessment}[] \xrightarrow{\text{Mediator}} \text{MediatorSynthesis} \\
    \text{R2:} \quad &\{\text{Self R1} + \text{MediatorSynthesis}\} \xrightarrow{\parallel} \text{RevisedAssessment}[] \\
    \text{Judge:} \quad &\text{RevisedAssessment}[] \xrightarrow{\text{aggregation}} \text{RankedPrediction}[]
\end{aligned}
\end{equation}

Каждая стадия:
\begin{itemize}
    \item Независима по данным (агенты не видят выводов друг друга до медиации).
    \item Парсируется в структурированный Pydantic-объект для валидации.
    \item Логируется с полной стоимостью LLM (tokens, USD).
    \item Переносит цепочку рассуждений, а не только вероятности.
\end{itemize}

Полная спецификация Delphi: `docs/05-delphi-pipeline.md`.



\newpage
% =====================================================================
% PART V: Aggregation & Generation (Stages 6-9)
% =====================================================================
% Part V: Агрегация и генерация (Stages 6-9)

\section{Judge: агрегация и ранжирование}
\label{sec:judge}

Стадия~6 пайплайна реализует финальную агрегацию персональных прогнозов~\emph{Delphi ensemble} и выбор событий для генерации заголовков. Двухшаговый алгоритм:

\subsection{Шаг 6a: Агрегация на уровне событий}

\subsubsection{Парсинг оценок}

Агент Judge принимает результаты раунда~2 (Round~2) либо раунда~1 (Round~1) в качестве fallback. Каждая оценка \verb|PersonaAssessment| содержит список прогнозов конкретной персоны с вероятностями.

\subsubsection{Взвешенная медиана}

Для каждого события (идентифицируемого \verb|event_thread_id|) агент вычисляет взвешенную медиану вероятностей пяти персон:
\begin{itemize}
  \item \textbf{REALIST}: базовый вес 0.22
  \item \textbf{GEOSTRATEG}: базовый вес 0.20
  \item \textbf{ECONOMIST}: базовый вес 0.20
  \item \textbf{MEDIA\_EXPERT}: базовый вес 0.18
  \item \textbf{DEVILS\_ADVOCATE}: базовый вес 0.20
\end{itemize}

\subsubsection{Horizon-адаптивные коррекции}

Веса корректируются в зависимости от временного горизонта (расстояния до целевой даты):

\begin{table}[h]
  \centering
  \caption{Horizon-адаптивные корректировки весов персон}
  \label{tab:horizon-weights}
  \begin{tabular}{|l|c|c|c|c|c|}
    \hline
    \textbf{Горизонт} & \textbf{REALIST} & \textbf{GEO} & \textbf{ECON} & \textbf{MEDIA} & \textbf{DEVIL} \\
    \hline
    \emph{Immediate} (0--2 дня) & $-0.02$ & $-0.02$ & $+0.02$ & $+0.03$ & $-0.01$ \\
    \hline
    \emph{Near} (3--4 дня) & $0.00$ & $0.00$ & $0.00$ & $0.00$ & $+0.02$ \\
    \hline
    \emph{Medium} (5+ дней) & $+0.03$ & $+0.02$ & $-0.01$ & $-0.02$ & $-0.02$ \\
    \hline
  \end{tabular}
\end{table}

Обоснование: для немедленных прогнозов важны циклы новостей (Media Expert) и календарные факторы (Economist); для дальних горизонтов -- структурные тренды (Realist, Geostrateg).

\subsubsection{Инъекция рыночного сигнала}

Если в контексте пайплайна присутствуют \emph{foresight signals} (Polymarket), агент выполняет нечёткое сопоставление события с рынками через метрику Jaccard:
\begin{itemize}
  \item Tier~1: точное совпадение (similarity $\geq 0.65$)
  \item Tier~2: частичное совпадение (similarity $\geq 0.40$ и Jaccard $\geq 0.3$)
\end{itemize}

Если совпадение найдено, добавляется псевдо-персона <<Market>> с динамическим весом:
\[
w_{\text{market}} = 0.15 \times \max\!\bigl(0.5,\; \min(\log_{10}(V)/6,\; 1.5)\bigr) \times \frac{1}{1 + \text{vol}} \times \text{reliable} \times r_{\text{horizon}}
\]

где $V$ --- объём рынка, $r_{\text{horizon}}$ учитывает дни до закрытия рынка, $\text{reliable} = 0$ если распределение ненадёжно, иначе $1$.

Market-сигнал инъектируется только при $p_{\text{market}} \geq 0.10$ (\texttt{MARKET\_MIN\_PROBABILITY}) для исключения шума от неликвидных рынков. Дополнительно применяется фильтр ликвидности: рынки с объёмом $< \$10{,}000$ (\texttt{MARKET\_MIN\_LIQUIDITY}) исключаются при построении индекса.

\subsubsection{Калибровка: Platt Scaling с экстремизацией}

После вычисления взвешенной медианы $p_{\text{raw}}$ применяется \emph{Platt scaling} с параметром экстремизации $a = 1.5$:
\begin{equation}
p_{\text{calibrated}} = \sigma(a \cdot \text{logit}(p_{\text{raw}}) + b)
\label{eq:platt}
\end{equation}

где $\text{logit}(p) = \ln(p/(1-p))$, $\sigma(x) = 1/(1+e^{-x})$, $b = 0.0$.

Экстремизация ($a > 1$) усиливает уверенные прогнозы и ослабляет неуверенные, повышая калибровку на реальных данных.

\subsubsection{Уровни согласия}

Разброс вероятностей $\text{spread} = \max(p_i) - \min(p_i)$ классифицируется как:
\begin{itemize}
  \item \textbf{CONSENSUS}: $\text{spread} < 0.15$ (тесное согласие)
  \item \textbf{MAJORITY\_WITH\_DISSENT}: $0.15 \leq \text{spread} < 0.30$
  \item \textbf{CONTESTED}: $\text{spread} \geq 0.30$ (глубокие разногласия)
\end{itemize}

При CONTESTED применяется штраф неопределённости ($\times 0.8$) к сырой вероятности.

\subsubsection{Агрегация временных полей}

Для каждого события агрегируются поля: \verb|predicted_date| (медиана по датам персон), \verb|uncertainty_days| (среднее), \verb|scenario_types| (объединение), \verb|causal_dependencies| (объединение). Выход: \verb|PredictedTimeline| с отсортированными по дате \verb|TimelineEntry|.

\subsection{Шаг 6b: Отбор заголовков}

\subsubsection{Оценка новостности}

Для каждого события в timeline вычисляется headline score:
\begin{equation}
\text{score} = p_{\text{calib}} \times n \times (1 - s) \times r
\label{eq:headline-score}
\end{equation}

где:
\begin{itemize}
  \item $p_{\text{calib}}$ — калиброванная вероятность события
  \item $n$ — новостная ценность (среднее по персонам)
  \item $s$ — коэффициент насыщенности (текущее заполнение 0.0)
  \item $r$ — релевантность для издания (по умолчанию 1.0)
\end{itemize}

\subsubsection{Временная близость}

Коэффициент temporal proximity штрафует события, далёкие от целевой даты:
\[
\text{temporal\_factor} = \max(0.5, 1.0 - 0.05 \cdot \text{days\_away})
\]

\subsubsection{Топ-7 + дикие карты}

Отбираются 7 событий с наивысшим финальным score. Дополнительно добавляются до 2 <<дикие карты>> (black swan сценарии от Devil's Advocate) с новостной ценностью $\geq 0.7$. Вывод: \verb|RankedPrediction|[] с полями \verb|headline_score|, \verb|rank|, \verb|is_wild_card|.

---

\section{FramingAnalyzer: редакционный угол}
\label{sec:framing}

Стадия~7 анализирует, \emph{как конкретное издание подаст событие}. Один и тот же факт получает разные редакционные углы в зависимости от политики и стиля издания.

\subsection{Девять стратегий фрейминга}

\begin{table}[h]
  \centering
  \caption{Стратегии фрейминга в FramingAnalyzer}
  \label{tab:framing-strategies}
  \begin{tabular}{|l|l|}
    \hline
    \textbf{Стратегия} & \textbf{Пример} \\
    \hline
    \emph{THREAT} & Угроза для здоровья, безопасности, экономики \\
    \hline
    \emph{OPPORTUNITY} & Шанс, перспектива развития \\
    \hline
    \emph{CRISIS} & Острая проблема, обострение ситуации \\
    \hline
    \emph{ROUTINE} & Регулярное, периодическое событие \\
    \hline
    \emph{SENSATION} & Шокирующие подробности, неожиданный поворот \\
    \hline
    \emph{ANALYTICAL} & Глубокий анализ, контекст, причины \\
    \hline
    \emph{HUMAN\_INTEREST} & История людей, личные драмы \\
    \hline
    \emph{NEUTRAL\_REPORT} & Фактический отчёт без оценки \\
    \hline
    \emph{CONFLICT} & Столкновение позиций, дебаты \\
    \hline
  \end{tabular}
\end{table}

\subsection{Структура FramingBrief}

Для каждого \verb|RankedPrediction| генерируется \verb|FramingBrief| с полями:

\begin{itemize}
  \item \textbf{framing\_strategy}: одна из девяти стратегий
  \item \textbf{angle}: конкретный угол подачи (1--2 предложения)
  \item \textbf{emphasis\_points}: 2--5 элементов, которые издание подчеркнёт
  \item \textbf{omission\_points}: 0--5 элементов, которые издание приглушит или опустит
  \item \textbf{headline\_tone}: тревожный / нейтральный / оптимистичный / ироничный
  \item \textbf{likely\_sources}: 1--5 источников, на которые сошлётся издание
  \item \textbf{section}: раздел издания (Политика, Экономика, Главное и т.д.)
  \item \textbf{news\_cycle\_hook}: привязка к текущему новостному циклу
  \item \textbf{editorial\_alignment\_score}: насколько событие соответствует редакционной линии (0 -- совсем не в профиле, 1 -- идеально)
\end{itemize}

\subsection{Входные данные}

Для анализа используется \verb|OutletProfile| с:
\begin{itemize}
  \item Редакционной позицией (editorial\_stance)
  \item Типичными разделами (sections)
  \item 10--20 примерами недавних заголовков (sample\_headlines)
  \item Тональным профилем (tone\_profile: нейтральный, сенсационный, аналитический)
  \item Типичными источниками (sourcing\_patterns)
\end{itemize}

Модель: Claude Sonnet (OpenRouter) с temperature 0.5 для консервативного reasoning.

---

\section{StyleReplicator: генерация заголовков}
\label{sec:style-replicator}

Стадия~8 генерирует 2--3 варианта заголовка и первого абзаца в стиле целевого издания.

\subsection{Few-shot инъекция стиля}

StyleReplicator получает из \verb|OutletProfile|:
\begin{itemize}
  \item 10--20 реальных заголовков издания (примеры для имитации)
  \item 3--5 примеров первых абзацев (лидов)
  \item Метрики стиля:
    \begin{itemize}
      \item avg\_headline\_length: средняя длина в символах
      \item headline\_case: нижний / верхний / sentence case / title case
      \item tone\_profile: нейтральный / острый / сенсационный / аналитический
      \item quotes\_usage: частота использования кавычек
      \item colon\_usage: частота двоеточия
      \item avg\_lede\_length: средняя длина первого абзаца в словах
    \end{itemize}
\end{itemize}

\subsection{Мультиязычность}

Выбор модели зависит от языка издания:
\begin{table}[h]
  \centering
  \caption{Выбор модели по языку издания}
  \label{tab:language-models}
  \begin{tabular}{|l|l|}
    \hline
    \textbf{Язык} & \textbf{Модель} \\
    \hline
    Русский & Claude Sonnet (OpenRouter) \\
    \hline
    Английский & Claude Sonnet (OpenRouter) \\
    \hline
    Другие & Claude Sonnet (OpenRouter) \\
    \hline
  \end{tabular}
\end{table}

\subsection{Промпт структура}

Системный промпт содержит идентичность (<<Ты -- автор заголовков для издания X>>), параметры стиля и примеры. Пользовательский промпт:
\begin{itemize}
  \item Событие и его вероятность
  \item FramingBrief (стратегия, угол, тон, источники)
  \item Запрос: сгенерировать 2--3 варианта с первым абзацем
\end{itemize}

Температура: 0.8 для творческого разнообразия. Выход: \verb|GeneratedHeadline|[] с полями \verb|headline|, \verb|first_paragraph|, \verb|headline_language|, \verb|length_deviation|.

---

\section{QualityGate: контроль качества}
\label{sec:quality-gate}

Стадия~9 — финальный фильтр. Три независимые проверки для каждого заголовка, принимающие решение: PASS / REJECT / REVISE / DEPRIORITIZE / MERGE.

\subsection{Проверка 1: Фактическая правдоподобность}

\subsubsection{Задача}

Оценить внутреннюю непротиворечивость заголовка и первого абзаца, соответствие известным фактам. \emph{НЕ} оценивает, произойдёт ли событие.

\subsubsection{Критерии}

\begin{itemize}
  \item Нет ли противоречий с общеизвестными фактами?
  \item Корректны ли упомянутые даты, должности, названия?
  \item Логически непротиворечив ли сценарий?
  \item Не является ли событие физически / логически невозможным?
  \item Не устарела ли предпосылка?
\end{itemize}

\subsubsection{Оценка (1--5)}

\begin{table}[h]
  \centering
  \caption{Шкала фактической проверки}
  \label{tab:factual-score}
  \begin{tabular}{|c|l|}
    \hline
    \textbf{Балл} & \textbf{Интерпретация} \\
    \hline
    5 & Полностью правдоподобно, фактических противоречий не найдено \\
    \hline
    4 & Мелкие неточности, не влияющие на смысл \\
    \hline
    3 & Спорные утверждения, но принципиально не ошибочные \\
    \hline
    2 & Фактическая ошибка или серьёзное логическое противоречие \\
    \hline
    1 & Явно абсурдный или невозможный сценарий \\
    \hline
  \end{tabular}
\end{table}

\textbf{Порог}: $\text{factual\_score} < 3$ $\to$ REJECT (проблема в прогнозе, а не в генерации).

Модель: Claude Sonnet, temperature 0.2 (консервативно).

\subsection{Проверка 2: Стилистическая аутентичность}

\subsubsection{Задача}

Оценить, насколько заголовок и первый абзац стилистически соответствуют целевому изданию.

\subsubsection{Критерии}

\begin{itemize}
  \item Типична ли лексика для этого издания?
  \item Соответствуют ли синтаксические конструкции?
  \item Подходит ли длина заголовка?
  \item Верный ли тон?
  \item <<Smell test>>: похоже ли это на реальную публикацию этого издания?
\end{itemize}

\subsubsection{Оценка (1--5)}

\begin{table}[h]
  \centering
  \caption{Шкала стилистической проверки}
  \label{tab:style-score}
  \begin{tabular}{|c|l|}
    \hline
    \textbf{Балл} & \textbf{Интерпретация} \\
    \hline
    5 & Неотличимо от настоящей публикации \\
    \hline
    4 & Стиль в целом верный, мелкие шероховатости \\
    \hline
    3 & Узнаваемый стиль, но с заметными отклонениями \\
    \hline
    2 & Стиль не соответствует изданию \\
    \hline
    1 & Совершенно не похоже на это издание \\
    \hline
  \end{tabular}
\end{table}

\textbf{Порог}: $\text{style\_score} < 3$ $\to$ REVISE (одна попытка ревизии в StyleReplicator).

Модель: Claude Sonnet (для всех языков), temperature 0.2.

\subsection{Проверка 3: Дедупликация}

\subsubsection{Внутренняя дедупликация}

Вычисляются embeddings всех заголовков (text-embedding-3-small через OpenRouter). Если cosine similarity между двумя заголовками $> 0.85$, слабее по среднему score помечается дубликатом.

\subsubsection{Внешняя дедупликация}

Если cosine similarity заголовка с реальными заголовками из \verb|OutletProfile.sample_headlines| $> 0.80$, заголовок помечается как external duplicate.

\subsection{Матрица решений}

\begin{table}[h]
  \centering
  \caption{Логика гейта QualityGate}
  \label{tab:gate-logic}
  \begin{tabular}{|c|c|c|c|}
    \hline
    \textbf{factual} & \textbf{style} & \textbf{dup} & \textbf{Решение} \\
    \hline
    $\geq 3$ & $\geq 3$ & нет & \textbf{PASS} \\
    \hline
    $< 3$ & любой & любой & \textbf{REJECT} \\
    \hline
    $\geq 3$ & $< 3$ & нет & \textbf{REVISE} \\
    \hline
    $\geq 3$ & $\geq 3$ & внутр. & \textbf{MERGE} \\
    \hline
    $\geq 3$ & $\geq 3$ & внеш. & \textbf{DEPRIORITIZE} \\
    \hline
  \end{tabular}
\end{table}

\subsection{Обработка решений}

\begin{itemize}
  \item \textbf{PASS}: заголовок добавляется в \verb|FinalPrediction|
  \item \textbf{REJECT}: заголовок отклоняется (без повторных попыток)
  \item \textbf{REVISE}: заголовок отправляется обратно в StyleReplicator с обратной связью (максимум 1 попытка), затем перепроверяется стиль
  \item \textbf{MERGE}: дубликат не добавляется, более сильный вариант уже в списке
  \item \textbf{DEPRIORITIZE}: заголовок добавляется в конец списка
\end{itemize}

\subsection{Выход}

На входе: 14--21 вариант заголовков (2--3 на каждое из 7 событий).
На выходе: 5--7 \verb|FinalPrediction| после фильтрации.

Каждый \verb|FinalPrediction| содержит:
\begin{itemize}
  \item Основной заголовок и первый абзац
  \item 1--2 альтернативных заголовка
  \item Калиброванную вероятность события
  \item Уровень согласия экспертов
  \item Стратегию фрейминга
  \item Язык заголовка
  \item Флаг wild card
\end{itemize}

Параллельная обработка: до 5 заголовков проверяются одновременно с timeout 180 сек. При timeout применяется нейтральная оценка (score = 3, пропускает заголовок дальше).

% end of Part V

\newpage
% =====================================================================
% PART VI: Polymarket Signal Enrichment (existing content, as-is)
% =====================================================================
\input{parts/part_vi}

\newpage
% =====================================================================
% PART VIII: Evaluation
% =====================================================================
\section{Brier Score: основная метрика}

Калибровка вероятностных предсказаний измеряется стандартной метрикой \emph{Brier Score} (BS), изначально разработанной для оценки метеорологических прогнозов.

\begin{equation}
BS = \frac{1}{N} \sum_{i=1}^{N} (f_i - o_i)^2
\label{eq:brier_score}
\end{equation}

где $f_i \in [0,1]$ "--- предсказанная вероятность события $i$, $o_i \in \{0,1\}$ "--- фактический исход (1, если событие произошло, 0 "--- в противном случае).

В контексте Delphi Press предсказываемым событием является наличие освещения прогнозируемой темы в целевом медиаиздании на заданный горизонт времени (1, 3 или 7 дней). Вероятность $f_i$ берётся из поля \texttt{FinalPrediction.confidence}, а бинаризация исхода $o_i$ определяется как:

\begin{equation}
o_i = \begin{cases}
1, & \text{если } \texttt{TopicMatch} > 0 \\
0, & \text{иначе}
\end{cases}
\end{equation}

Источники для ground truth "--- актуальные RSS-потоки уже имеющихся 16 источников проекта (ТАСС RU/EN, РИА Новости, Интерфакс, BBC, Al Jazeera, Guardian и др.) и Wayback Machine CDX API, позволяющий восстанавливать снимки веб-страниц за прошлые даты.

Для оценки навыка система использует \emph{Brier Skill Score}:

\begin{equation}
BSS = 1 - \frac{BS}{0.25}
\end{equation}

где 0.25 "--- базовая ошибка случайного угадывания (всегда вероятность 50\%). BSS интерпретируется как относительное улучшение: $BSS = 0$ означает производительность на уровне случайности, $BSS = 1$ "--- идеальное предсказание.

Таблица~\ref{tab:brier_benchmarks} приводит калибровочные ориентиры из актуальных источников (состояние на апрель 2026 года):

\begin{table}[h!]
\centering
\caption{Benchmarks Brier Score для прогнозных систем}
\label{tab:brier_benchmarks}
\begin{tabular}{lrrl}
\toprule
\textbf{Система/уровень} & \textbf{Brier Score} & \textbf{BSS} & \textbf{Контекст} \\
\midrule
Случайное угадывание & 0.250 & 0.00 & Базовая линия (всегда 50\%) \\
Polymarket (prediction market) & $\sim$0.187 & $\sim$0.25 & Реальные деньги, мегаполис участников\footnote{\citep{clinton2024polymarket}} \\
Metaculus-сообщество (AI-вопросы) & 0.182 & $\sim$0.27 & Краудсорсированные прогнозы [self-reported] \\
Good Judgment Project (год 2) & $\sim$0.140 & $\sim$0.44 & Обученные суперпрогнозисты\footnote{\citep{tetlock2015superforecasting}} \\
\textbf{Delphi Press v1.0 (цель)} & \textbf{$<$0.20} & \textbf{$>$0.20} & Уровень prediction market \\
GPT-4.5 (ForecastBench, окт 2025) & 0.101 & $\sim$0.60 & Лучший LLM на сегодня [ForecastBench] \\
\textbf{Delphi Press v2.0 (цель)} & \textbf{$<$0.12} & \textbf{$>$0.52} & Уровень обученных прогнозистов \\
Суперпрогнозисты (ForecastBench) & 0.068--0.086 & 0.66--0.73 & Долгосрочный ориентир\footnote{\citep{tetlock2015superforecasting}; числа из ForecastBench} \\
\bottomrule
\end{tabular}
\end{table}

\textbf{Статус внедрения:} Пилотная оценка требует сбора ground truth, что предполагает автоматизацию запроса к Wayback Machine CDX API и парсинга RSS-снимков. Пилотный тест спланирован на объёме 50 prediction runs по 3 горизонтам, что даст примерно 150--350 пар событие-исход. Статистическая мощность при N=50 обеспечивает 95\%-й доверительный интервал ±0.03--0.05 вокруг BS. \textbf{Стоимость пилота: менее \$1} (только LLM-вызовы для граничных случаев TopicMatch; BERTScore и Wayback API бесплатны).


\section{Композитная оценка}

Так как Brier Score оценивает лишь калибровку вероятностей, а не качество самого предсказанного текста, для полной оценки качества введена взвешенная композитная метрика:

\begin{equation}
\text{CompositeScore} = 0.40 \times \text{TopicMatch} + 0.35 \times \text{SemanticSim} + 0.25 \times \text{StyleMatch}
\label{eq:composite}
\end{equation}

Компоненты метрики описаны в таблице~\ref{tab:composite_components}:

\begin{table}[h!]
\centering
\caption{Компоненты CompositeScore}
\label{tab:composite_components}
\begin{tabular}{llll}
\toprule
\textbf{Компонент} & \textbf{Диапазон} & \textbf{Вычисление} & \textbf{Вес} \\
\midrule
TopicMatch & $\{0.0, 0.5, 1.0\}$ & Тема совпадает с реальностью & 0.40 \\
SemanticSim & $[0.0, 1.0]$ & BERTScore F1 vs реальный заголовок & 0.35 \\
StyleMatch & $[0.0, 1.0]$ & LLM-as-judge оценка стилевого соответствия & 0.25 \\
\bottomrule
\end{tabular}
\end{table}

\textbf{TopicMatch} "--- трёхступенчатый алгоритм с возрастающей стоимостью:

\begin{enumerate}
    \item \textbf{Keyword screening} (< 5 мс): извлечение ключевых слов из предсказанного заголовка, фильтрация кандидатов из ground truth.
    \item \textbf{BERTScore} ($\sim$50 мс, локально): вычисление контекстного сходства через многоязычную модель \texttt{xlm-roberta-base} (для кроссъязычных пар).
    \item \textbf{LLM-арбитр} (только для пограничных случаев $0.55 \leq F1 < 0.60$): ~10\% случаев, что экономит ~90\% LLM-вызовов.
\end{enumerate}

Пороговые значения:
\begin{itemize}
    \item $F1 \geq 0.78 \Rightarrow$ TopicMatch = 1.0 (полное совпадение)
    \item $0.60 \leq F1 < 0.78 \Rightarrow$ TopicMatch = 0.5 (верная тема, но неточный фрейминг)
    \item $0.55 \leq F1 < 0.60 \Rightarrow$ LLM-арбитр (пограничные случаи)
    \item $F1 < 0.55 \Rightarrow$ TopicMatch = 0.0 (промах)
\end{itemize}

\textbf{SemanticSim} "--- F1-метрика BERTScore против лучшего совпадающего реального заголовка в день предсказания. Для монолингвальной оценки используются языкоспецифичные модели: \texttt{DeepPavlov/rubert-base-cased} для русского, \texttt{roberta-large} для английского.

\textbf{StyleMatch} "--- оценка по шкале 1--5 (затем нормализуется на [0,1] делением на 5) через LLM-as-judge. Критически: для избежания self-enhancement bias используется \textbf{отличная модель} от генератора заголовков. В стеке Delphi Press это Claude Sonnet 4.6 в роли судьи (в то время как Stage 8 Generator может использовать другую модель).

Интерпретация CompositeScore:
\begin{itemize}
    \item $\geq 0.70$ "--- отличный прогноз
    \item $0.50$--$0.69$ "--- хороший
    \item $0.30$--$0.49$ "--- частичное попадание
    \item $< 0.30$ "--- промах
\end{itemize}

\textbf{Известные ограничения:} Пороги BERTScore (0.78, 0.60) откалиброваны на задачах перевода и суммаризации, но не валидированы специально для русскоязычных новостных заголовков длиной 5--15 слов. Рекомендуется ручная аннотация 20--30 примеров в начале пилота. Покрытие Wayback Machine неоднородно: ТАСС и РИА Новости имеют хорошее архивное покрытие, Коммерсант и Ведомости "--- частичное.


\section{Walk-forward валидация inverse problem}

Для доказательства того, что информированный консенсус бетторов Polymarket действительно даёт сигнал, применяется протокол walk-forward валидации с непересекающимися окнами. Это предотвращает look-ahead bias, при котором модель "видит" будущие данные.

\textbf{Примечание о метрике.} В отличие от теоретического BSS из Раздела~1 (baseline = random 0.25), walk-forward BSS использует raw market price как baseline: $\bss_{wf} = 1 - \bs(\text{informed}) / \bs(\text{raw market})$.

\textbf{Протокол:}
\begin{itemize}
    \item \textbf{Burn-in:} 180 дней исторических данных для обучения профилей (безопасный зазор перед тестом).
    \item \textbf{Фолды:} 22 непересекающихся 60-дневных тестовых окна, с шагом 60 дней.
    \item \textbf{Данные:} 435K разрешённых рынков, 470M торговых операций, сведённых в 82.6M временных агрегатов в parquet-файле с 30-дневными бакетами.
    \item \textbf{Временная утечка:} Предыдущие реализации использовали предварительно агрегированные позиции, включающие торги, совершённые \textbf{после} временной точки отсечки T. Исправление: пересканирование 33 ГБ raw trades в 30-дневные частичные суммы, из которых позиции at-cutoff восстанавливаются SQL-фильтром \texttt{WHERE time\_bucket <= T/bucket\_size}.
\end{itemize}

\textbf{Инфраструктура:}

Выбран DuckDB с потоковой обработкой и predicate pushdown на \texttt{time\_bucket} для избежания загрузки 470M торговых операций в оперативную память. Bucket-подход даёт 225$\times$ ускорение (15 мин на фолд $\to$ 4 сек) и снижает пик памяти с 7.4 ГБ (crash at fold 15) до 4.6 ГБ.

Результаты валидации представлены в таблице~\ref{tab:walkforward_bss}:

\begin{table}[h!]
\centering
\caption{Walk-forward BSS: 22 фолда с временными бакетами}
\label{tab:walkforward_bss}
\begin{tabular}{lrrrr}
\toprule
\textbf{Подмножество} & \textbf{Mean BSS} & \textbf{Median BSS} & \textbf{Fraction > 0} & \textbf{95\% CI} \\
\midrule
All 22 folds & +0.196 & +0.159 & 22/22 (100\%) & [+0.094, +0.297] \\
Robust (folds 0--16, test $\geq$944) & +0.127 & +0.095 & 17/17 (100\%) & — \\
\bottomrule
\end{tabular}
\end{table}

\textbf{Интерпретация:}

\begin{enumerate}
    \item \textbf{100\% положительных фолдов (22/22).} Информированный консенсус \textbf{всегда} улучшает сырую рыночную цену Polymarket, без исключений.

    \item \textbf{Robust mean BSS = +0.127.} На подмножестве из 17 фолдов с надёжным объёмом тестовых рынков ($\geq 944$) улучшение Brier Score составляет 12.7\% относительно сырых цен.

    \item \textbf{Инвертированная U-кривая.} BSS достигает пика при fold 9--11 (+0.21--+0.27) с оптимальным объёмом данных (~35K информированных бетторов), затем уменьшается по мере увеличения ликвидности рынка. Это соответствует литературе (Akey et al. 2025): информированные торговцы более значимы на тонких рынках.

    \item \textbf{Стабильность яруса.} Метрика Jaccard overlap между множествами INFORMED бетторов в соседних фолдах = 0.613 (61\%). Превышает пороговое значение 60\%, подтверждая стабильность классификации яруса.

    \item \textbf{Масштаб.} 1.74M профилированных кошельков, из них 348K в ярусе INFORMED "--- крупнейший published dataset бетторских профилей на Polymarket.
\end{enumerate}

\textbf{Варианты без extremizing.} Абляция (Phase 5 в CHANGELOG) показала, что простая accuracy-weighted consensus с Bayesian shrinkage (k=15) оптимальна. Экстремизация (Satopää et al. 2014) фактически \textbf{вредит} (-64\% BSS), вероятно потому, что информированные бетторы коррелированы (реагируют на те же рыночные сигналы), а не независимо информированы (частная информация).

\textbf{Корреляция новостей и рынка.}

Дополнительная валидация использует Spearman rank correlation и Granger causality test для проверки направления влияния новостного сигнала на цены:

\begin{itemize}
    \item \textbf{Spearman $\rho$:} между абсолютным изменением цены $|\Delta p|$ и количеством релевантных новостей (требует min 5 datapoints).
    \item \textbf{Granger causality:} тестирует, предсказывает ли объём новостей изменения цены с лагом до 7 дней. ADF test проверяет стационарность; результат "--- (F-статистика, p-value, оптимальный лаг) или (None, None, None) при недостаточных данных.
\end{itemize}

Эти метрики вычисляются в \texttt{src/eval/correlation.py} и используются для диагностики, но не входят в основную Brier Score оценку.


\newpage
% =====================================================================
% PART IX: Dead Ends & Lessons Learned
% =====================================================================
\section{API dead ends}

\subsection{Case Study 1: Metaculus 403 API Deprecation and Tier Lock (v0.5.1)}

\subsubsection*{Контекст}
Metaculus служит источником структурированных прогнозов сообщества в stage 1 (ForesightCollector). Их API интегрирована в \texttt{src/data\_sources/foresight.py} с 2026-03-28. На production сервере v0.5.1 все запросы к Metaculus возвращали HTTP 403 Forbidden, но локально тесты проходили. Это свидетельствовало о скрытом breaking change в API.

\subsubsection*{Что пробовали}
В диагностическом скрипте (\texttt{tasks/research/metaculus\_api\_fix.md}) проверили несколько гипотез:
\begin{enumerate}
\item GET requests требуют авторизацию (было предположение: read-доступ свободный)
\item Параметры изменили названия между версиями (`status` vs `statuses`)
\item Endpoint \texttt{/api2/questions/} всё ещё работает как документировано
\end{enumerate}

\subsubsection*{Почему не сработало}
Две причины одновременно:

\textbf{(1) API v2 fully deprecated.} Примерно в Q3--Q4 2024 Metaculus перенёсли все endpoints с \texttt{/api2/*} на \texttt{/api/posts/*}. Документация не синхронизирована. Endpoint \texttt{/api2/questions/} остался работающим 10--15 месяцев, создав впечатление обратной совместимости. На самом деле, поддержка заканчивается.

\textbf{(2) Все endpoints требуют авторизацию.} Историческая норма: GET-запросы к публичным API бесплатные без токена. Metaculus изменили политику --- даже чтение вопросов требует `Authorization: Token <token>` заголовок. Возможно, это было сделано для rate-limiting по пользователю.

\subsubsection*{Что сделали}
Fix реализован в v0.5.1 (коммит в CHANGELOG.md, стр. 15):

\begin{enumerate}
\item Конструктор MetaculusClient принимает опциональный параметр \texttt{token: str = ""}. Если указан, заголовок \texttt{Authorization: Token {token}} добавляется в каждый запрос.
\item Endpoint изменён с \texttt{/api2/questions/} на \texttt{/api/posts/}.
\item Параметры переименованы:
    \begin{itemize}
    \item `status` $\to$ `statuses`
    \item \texttt{resolve\_time\_\_gt} $\to$ \texttt{scheduled\_resolve\_time\_\_gt}
    \item \texttt{resolve\_time\_\_lt} $\to$ \texttt{scheduled\_resolve\_time\_\_lt}
    \end{itemize}
\item Парсинг ответа переписан. Новый формат:
\begin{verbatim}
post['question']['aggregations']['recency_weighted']['latest']['centers'][0]
\end{verbatim}
вместо старого:
\begin{verbatim}
post['community_prediction']['full']['q2']
\end{verbatim}
\item Токен генерируется на https://www.metaculus.com/aib/ (бесплатный, не истекает).
\end{enumerate}

v0.9.4 (CHANGELOG, стр. 83) Metaculus отключён в production (\texttt{\_fetch\_metaculus()} возвращает \texttt{[]} если client=None). Причина: access tier ограничен --- нам нужен BENCHMARKING tier для надёжного доступа.

\subsubsection*{Урок}
\textbf{Никогда не предполагай обратную совместимость третьего-сторонней API.} Даже если старый endpoint работает 10+ месяцев, это может быть legacy layer, а не guarantee.
\begin{enumerate}
\item Подписывайся на изменения API через официальные каналы (RSS, email newsletter, GitHub releases).
\item При интеграции внешней API, сохраняй последнюю документацию локально (\texttt{tasks/research/metaculus\_api\_fix.md} служит reference).
\item Зафиксируй rate limit --- Metaculus: 120 req/min (безопасно), без опубликованного лимита. Всегда консервативен.
\item При auth-требовании, явно документируй, как получить credentials, и сделай их конфигурируемыми из окружения.
\end{enumerate}

---

\subsection{Case Study 2: GDELT Cyrillic Query Crash (pre-v0.5.1, v0.9.4)}

\subsubsection*{Контекст}
ForesightCollector вызывает GDELT DOC 2.0 API с query вида \emph{``news forecast 2026-03-29 ТАСС''} для сбора статей по русским СМИ. На production сервере с v0.9.2 и v0.9.3 это вызывало JSON parse error: ElasticSearch backend за GDELT возвращал HTML вместо JSON.

\subsubsection*{Что пробовали}
В \texttt{tasks/research/gdelt\_api\_fix.md} и \texttt{tasks/research/gdelt\_api.md} проведён анализ:
\begin{enumerate}
\item GDELT query-язык поддерживает Unicode кириллицу?
\item API имеет опции для фильтрации по языку источника?
\item Как парсить 403 HTML-ответ, чтобы не падать?
\item Есть ли более структурированный источник событий (не pure search API)?
\end{enumerate}

\subsubsection*{Почему не сработало}
Три причины одновременно:

\textbf{(1) Кириллица в query.} GDELT индексирует только английские переводы и метаданные. Когда ForesightCollector отправляет \texttt{query="news forecast 2026-03-29 ТАСС"}, ElasticSearch не может обработать кириллицу в текущей конфигурации. Вместо JSON error, API возвращает HTML-страницу (5xx error page).

\textbf{(2) Будущие даты бессмысленны.} GDELT основан на rolling window ~ 3 месяца в прошлое. Query с датой 2026-03-29 (из futuer) не вернёт никаких статей. API не валидирует дату --- просто возвращает пустой результат (или ошибку).

\textbf{(3) Отсутствие защиты от null в `articles`.} Когда API испытывает нагрузку, может вернуть \texttt{{"articles": null}} вместо \texttt{{"articles": []}}. Код \texttt{for article in data.get("articles", [])} упадёт с TypeError, потому что `None` не iterable.

\subsubsection*{Что сделали}
Fix реализован в два этапа: Content-Type guard + null-safe articles в pre-v0.5.1 (commit 987260c), Cyrillic stripping в v0.9.4 (commit 3658531):

\begin{enumerate}
\item \textbf{Content-Type проверка.} Перед парсингом JSON:
\begin{verbatim}
if "text/html" in response.headers.get("content-type", ""):
    logger.warning("GDELT returned HTML: %s", response.text[:200])
    return []
\end{verbatim}

\item \textbf{Null-safe articles.} Замена:
\begin{verbatim}
for article in (data.get("articles") or []):
\end{verbatim}
вместо \texttt{data.get("articles", [])}, потому что latter не защищает от `null`.

\item \textbf{Английский query.} ForesightCollector теперь генерирует query на английском: `"news forecast"` вместо `"Новостной прогноз"`. Language-specific фильтрация делается через GDELT operators (ниже).

\item \textbf{Language и country operators.} Документация в \texttt{tasks/research/gdelt\_api.md} (стр. 8--9) показывает:
\begin{verbatim}
sourcelang:russian + sourcecountry:RS
\end{verbatim}
(RS = FIPS код России, не ISO RU). Это фильтрует статьи по языку источника и стране издателя без кириллицы в query.
\end{enumerate}

\subsubsection*{Урок}
\textbf{Никогда не доверяй charset-поддержке без явного тестирования.} Даже если API документация не упоминает ограничения, всегда проверяй граничные случаи (non-ASCII characters, future dates, null values).

\begin{enumerate}
\item Запрашиваемые языки обрабатывай через фильтр (e.g. `sourcelang:` operator), а не через query text.
\item Проверяй `content-type` заголовок перед парсингом JSON --- это спасает от HTML error pages.
\item Защищайся от null values: \texttt{(data.get("field") or [])} вместо \texttt{data.get("field", [])}.
\item Rate limit на GDELT: ~1 req/sec, без опубликованного лимита. Добавь exponential backoff на 429.
\end{enumerate}

---

\subsection{Case Study 3: Polymarket camelCase Parameter Mismatch (pre-v0.5.1)}

\subsubsection*{Контекст}
Polymarket Gamma API используется в ForesightCollector для fetch markets. Код в \texttt{src/data\_sources/foresight.py} (строка 165) отправлял:
\begin{verbatim}
GET /markets?order=volume_24hr
\end{verbatim}

На production это возвращало HTTP 422 Unprocessable Entity. Локально в dry-run не воспроизводилось (возможно, из-за кэша или другого контекста).

\subsubsection*{Что пробовали}
Live-запросы (2026-03-28) с разными вариантами параметра:
\begin{enumerate}
\item \texttt{order=volume\_24hr} (snake\_case, как документировано в algunos examples) $\to$ 422
\item \texttt{order=volume24hr} (camelCase, как в других частях Polymarket API) $\to$ 200
\item \texttt{order=volume} (generic, не специфично) $\to$ 200
\end{enumerate}

Вывод: параметр требует camelCase, но никаких доков не было.

\subsubsection*{Почему не сработало}
\textbf{Inconsistency между endpoints.} Polymarket имеет несколько API endpoints:
\begin{itemize}
\item \texttt{/markets} --- принимает camelCase (\texttt{volume24hr}, \texttt{liquidityNum}, \texttt{endDate})
\item \texttt{/events} --- принимает snake\_case (\texttt{volume\_24h}, \texttt{liquidity\_num})
\end{itemize}

Это расхождение нигде не задокументировано. Разработчик (справедливо) предположил, что snake\_case универсален. На самом деле, \texttt{/markets} настаивает на camelCase.

\subsubsection*{Что сделали}
Fix реализован в pre-v0.5.1 (commit 987260c, 2026-03-28):

Изменение в \texttt{src/data\_sources/foresight.py}:
\begin{verbatim}
# БЫЛО:
"order": "volume_24hr"

# СТАЛО:
"order": "volume24hr"
\end{verbatim}

Документация в \texttt{tasks/research/polymarket\_gamma\_api\_fix.md} (стр. 26--46) содержит полный reference валидных параметров для \texttt{/markets} endpoint.

\subsubsection*{Урок}
\textbf{API inconsistency --- это bug в дизайне API, не вашего кода.} Однако, когда ты интегрируешь, нужно быть готов к такому.

\begin{enumerate}
\item Всегда тестируй все endpoints той же API live перед production deploy.
\item Не полагайся на consistency между похожими endpoints --- каждый может иметь свою convention.
\item Создай локальный reference документ для каждого внешнего API (как \texttt{tasks/research/}) с примерами.
\item При 422 Unprocessable Entity, форматируй запрос как JSON и проверь каждый параметр через документацию.
\end{enumerate}

---

\subsection{Case Study 4: Reuters RSS Feed Deprecation (v0.9.4)}

\subsubsection*{Контекст}
Исторически GLOBAL\_RSS\_FEEDS включал Reuters фиды:
\begin{verbatim}
https://feeds.reuters.com/reuters/...
\end{verbatim}

На production после 2026-03-28 эти URLs возвращали 404, не обновляясь. Pipeline падал на stage 1 при попытке fetch news signals через эти dead feeds.

\subsubsection*{Что пробовали}
Проверили несколько гипотез:
\begin{enumerate}
\item Reuters сменили хост (e.g. reuters.com/news vs feeds.reuters.com)?
\item Feed требует авторизацию?
\item Хост DNS разрешается, но HTTP 404 возвращается?
\end{enumerate}

Итог: \texttt{feeds.reuters.com} закрыт полностью (HTTPS SSL certificate expired, не поддерживается). Сервис закрыт с 2020 года.

\subsubsection*{Почему не сработало}
Reuters мигрировала на внутренний ReutersConnect API (требует авторизация, недоступна для public). Public RSS feeds больше не поддерживаются.

\subsubsection*{Что сделали}
Удалены все Reuters feeds из \texttt{src/agents/collectors/news\_scout.py} (GLOBAL\_RSS\_FEEDS). Альтернативные источники (BBC, AP, Bloomberg) остаются.

\subsubsection*{Урок}
\textbf{RSS feeds --- живые ресурсы. Периодически проверяй.} Крупные издатели могут закрыть public feeds в любой момент.

\begin{enumerate}
\item Добавь health-check для RSS feeds (попытайся fetch за последнюю неделю).
\item Имей fallback источники --- Reuters закрыт, но BBC, AP, Bloomberg, свежие.
\item Не полагайся на одного издателя для news signals.
\item Отслеживай 404/410 в логах --- это сигнал, что feed мёртв.
\end{enumerate}

---

\section{Архитектурные отказы}

\subsection{Case Study 5: YandexGPT Stub (v0.8.0)}

\subsubsection*{Контекст}
v0.7.0 содержал условную логику для выбора LLM провайдера: если \texttt{OPENROUTER\_API\_KEY} отсутствует, fallback на YandexGPT (\texttt{ya\_gpt\_pro} или \texttt{ya\_gpt\_lite}). YandexGPT требует российский IP (или VPN) и YandexCloud аккаунт.

\subsubsection*{Почему не сработало}
\begin{enumerate}
\item Production сервер в Yandex Cloud (213.165.220.144) имеет доступ к YandexGPT, но интеграция бросала \texttt{NotImplementedError} при каждом вызове (никогда не была реализована дальше stub'а).
\item Приоритизация провайдера была жёсткая: если YandexGPT задан, OpenRouter не используется. Это создавало single point of failure.
\item Документация prompts отсутствовала специфична для YandexGPT (модель имеет своё поведение JSON parsing).
\item Обслуживание двух параллельных стеков (OpenRouter + YandexGPT) удваивает complexity.
\end{enumerate}

\subsubsection*{Что сделали}
v0.8.0 (CHANGELOG, стр. 231): YandexGPT полностью удален. Три задачи, которые раньше использовали YandexGPT (\texttt{style\_generation}, \texttt{style\_generation\_ru}, \texttt{quality\_style}), переведены на OpenRouter с \texttt{anthropic/claude-sonnet-4}.

\subsubsection*{Урок}
\textbf{Не строй архитектуру на наличии двух провайдеров, если один из них не работает.} Single LLM provider (OpenRouter) лучше, чем fallback на неработающего second.

\begin{enumerate}
\item Выбери один primary провайдер и придерживайся.
\item Если нужна redundancy, выбери два провайдера, оба fully tested и integrated.
\item Не добавляй условную логику ``если API key отсутствует, используй другого провайдера'' --- это скрывает problems.
\end{enumerate}

---

\subsection{Case Study 6: Non-Existent Preset Sonnet 4.6 (v0.9.4)}

\subsubsection*{Контекст}
v0.9.3 содержал 3 пресета в \texttt{src/config.py}:
\begin{enumerate}
\item Light (Gemini Flash)
\item Standard (Claude Sonnet 4.6)
\item Opus (Claude Opus 4.6)
\end{enumerate}

Standard пресет был фасадом для модели \texttt{claude-sonnet-4.6}, которой нет в OpenRouter pricing. Результат: при выборе Standard, pipeline падал на первом LLM call с ошибкой ``model not found''.

\subsubsection*{Почему не сработало}
\begin{enumerate}
\item Confusion между версиями Sonnet. OpenRouter поддерживает \texttt{claude-3.5-sonnet}, но не \texttt{claude-sonnet-4.6} (и не существует в природе).
\item Пресеты не были протестированы полностью (dry-run использовал флаги, не UI-выбор).
\item Документация была неактуальна (старые имена моделей).
\end{enumerate}

\subsubsection*{Что сделали}
v0.9.4 (CHANGELOG, стр. 85): Standard пресет удален. Оставлены 2 пресета: Light (Gemini 2.5 Flash) и Opus (Claude Opus 4.6). UI обновлена.

\subsubsection*{Урок}
\textbf{Валидируй model names до деплоя.} Если пресет ссылается на модель, проверь, что она существует в pricing.

\begin{enumerate}
\item Добавь CI/CD step, который тестирует каждый пресет с real API call.
\item Держи актуальный список доступных моделей на каждом провайдере.
\item Лучше иметь 1 пресет, который работает, чем 3, из которых 2 сломаны.
\end{enumerate}

---

\subsection{Case Study 7: Dark Mode Complexity (v0.8.0)}

\subsubsection*{Контекст}
v0.7.0 и ранее содержали CSS для dark mode (Pico.css имеет встроенную поддержку). Было HTML-переключение `<button>` для toggle. Это добавило:
\begin{enumerate}
\item CSS сложность (duplication для light/dark палеток)
\item JS сложность (localStorage для preference, system preference detection)
\item Testing complexity (тестировать оба режима)
\item User confusion (не все пользователи знают, что toggle существует)
\end{enumerate}

\subsubsection*{Почему не сработало}
\begin{enumerate}
\item Нет пользовательского спроса на dark mode (feedback не указывал на это).
\item Дизайн был субоптимален --- цвета светлого режима были тщательно выбраны (OKLCH, контрастность, доступность). Dark mode версия была скопирована без адаптации.
\item Maintenance burden: каждое изменение цвета требовало обновления двух палеток.
\end{enumerate}

\subsubsection*{Что сделали}
v0.8.0 (CHANGELOG, стр. 231): Dark mode удален. Оставлена только светлая тема (OKLCH палетка, оптимизированная для контрастности и доступности).

\subsubsection*{Урок}
\textbf{Не добавляй features без спроса. YAGNI.} Даже если feature ``nice to have'', это не стоит complexity.

\begin{enumerate}
\item Собирай feedback перед добавлением UI features.
\item Если feature не используется, удали её.
\item Лучше одна хорошо-сделанная тема, чем две плохо-сделанные.
\end{enumerate}

---

\subsection{Case Study 8: Миграция Pico.css на Tailwind CSS (v0.8.0)}

\subsubsection*{Контекст}
Начальная версия UI использовала Pico.css --- classless CSS-фреймворк (минимальная стилизация без классов). Для прототипов этого было достаточно: формы, таблицы и типографика выглядели приемлемо из коробки.

\subsubsection*{Почему не сработало}
С развитием продукта требования к UI выросли за пределы возможностей classless-фреймворка:
\begin{enumerate}
\item \textbf{OKLCH-палитра.} Pico.css работает на HSL/hex, не поддерживает OKLCH цветовое пространство, необходимое для perceptually uniform палитры дизайн-системы.
\item \textbf{Кастомные компоненты.} Dashboard с прогнозами, карточки результатов, timeline --- ни один из этих элементов не покрывается classless CSS.
\item \textbf{Responsive layout.} Pico.css предоставляет базовый responsive через media queries, но нет утилит для grid/flex раскладки dashboard'ов.
\item \textbf{Анимации.} Переходы между состояниями (loading, success, error) требуют кастомного CSS, что полностью нивелирует преимущество ``zero-class'' подхода.
\end{enumerate}

\subsubsection*{Что сделали}
v0.8.0 (CHANGELOG, стр. 231): полная миграция на Tailwind CSS v4 с PostCSS build pipeline (commit ae12705):

\begin{enumerate}
\item Tailwind CSS v4.2.2 с \texttt{@theme} config заменил Pico.css.
\item Реализована дизайн-система Impeccable с 17 JS-referenced \texttt{fn-*} компонентами.
\item Типографика: Newsreader (заголовки) / Source Sans 3 (body) / JetBrains Mono (код) через Google Fonts.
\item PostCSS build: \texttt{input.css} $\to$ \texttt{tailwind.css} с watch mode для разработки.
\end{enumerate}

\subsubsection*{Урок}
\textbf{Classless CSS-фреймворки --- ловушка для всего, что сложнее landing page.} Если UI потребует кастомных компонентов, утилитарный фреймворк нужно выбирать сразу.

\begin{enumerate}
\item Оценивай потолок UI-требований на старте проекта, а не по текущим потребностям.
\item Стоимость миграции CSS растёт нелинейно с количеством шаблонов.
\item Утилитарный фреймворк (Tailwind) масштабируется с проектом; classless --- нет.
\end{enumerate}

---

\section{Критические баги}

\subsection{Case Study 9: Temporal Leak in Walk-Forward Evaluation (v0.9.2)}

\subsubsection*{Контекст}
Phase 4 (v0.9.2) реализовала walk-forward evaluation для валидации informed consensus. Модель использовала pre-aggregated позиции бетторов (из старого JSON-датасета) за каждый рынок. При итерации по 22 фолдам walk-forward каждый фолд имеет temporal cutoff $T$ (``используй только данные до даты $T$'').

Проблема: pre-aggregated позиции содержали трейды со средней датой $T + 30$ дней. Это means информационная утечка из будущего --- модель видит сигналы, которые не могла видеть реальная система в момент времени $T$.

\subsubsection*{Что пробовали}
\begin{enumerate}
\item Проверить, что cutoff $T$ корректен? Да, cutoff применялась правильно.
\item Проверить, что фолды не перекрываются? Да, non-overlapping 60-day windows.
\item Где источник утечки? Traced в pre-aggregated \texttt{avg\_position} для каждого bettor-market пары.
\end{enumerate}

\subsubsection*{Почему не сработало}
\textbf{Design issue.} Pre-aggregated positions были computed один раз на весь датасет:
\begin{verbatim}
avg_position = SUM(price $\times$ quantity) / SUM(quantity) for all trades
\end{verbatim}

Это глобальное среднее, без временного измерения. При walk-forward, код использовал то же самое среднее для всех фолдов, забывая, что медиана датирована будущим.

\subsubsection*{Что сделали}
v0.9.2 (CHANGELOG, стр. 119--122): Новый скрипт \texttt{scripts/duckdb\_build\_bucketed.py} переагрегирует 470M raw trades в bucketed parquet:

\begin{enumerate}
\item Временные бакеты: 30-дневные окна (e.g., 2024-01-01 до 2024-01-31).
\item Для каждого (bettor, market, time\_bucket): сохрани \texttt{weighted\_price\_sum} и \texttt{total\_usd} (частичные суммы).
\item Walk-forward теперь может вычислить: $\text{avg\_position\_as\_of\_T} = \frac{\sum \text{weighted\_price\_sum}}{\sum \text{total\_usd}}$ где $\text{time\_bucket} \le T$.
\item Zero data loss: 470M trades $\to$ 2.4GB bucketed parquet с predicate pushdown.
\end{enumerate}

Result (CHANGELOG, стр. 127): 22/22 фолда $\text{BSS} > 0$ с robust mean $\text{BSS} = +0.127$.

Анализ утечки (CHANGELOG, стр. 130): leaked $\text{BSS} = +0.092$ vs clean $\text{BSS} = +0.117$ на тех же фолдах. Leak добавлял шум, не сигнал.

\subsubsection*{Урок}
\textbf{Temporal cutoff требует временного измерения в данных, не глобальное агрегирование.} Это критично для любого walk-forward evaluation.

\begin{enumerate}
\item Никогда не используй pre-aggregated data для temporal validation.
\item Bucketed/timestamped aggregates позволяют point-in-time queries: ``какое было среднее на 2024-01-15?''
\item Всегда проверяй: есть ли trades/signals с датой > cutoff T? Если да, это leak.
\item Unit test: add explicit check that \texttt{as\_of\_date} < cutoff date for all test data.
\end{enumerate}

---

\subsection{Case Study 10: conditionId Mismatch (v0.9.3)}

\subsubsection*{Контекст}
Polymarket имеет дуальную структуру: каждый рынок имеет two IDs:
\begin{itemize}
\item `id` (numeric, e.g., 123456) --- локальный в Gamma API marketplace
\item `conditionId` (CTF hex hash, e.g., 0x7e2a...) --- глобальный в CLOB и Data API
\end{itemize}

Phase 6 (v0.9.5, но bug был в v0.9.3) интегрировала Data API для fetch real trades. ForesightCollector пытается обогатить Gamma markets с informèd consensus signals из trade data.

Проблема: \texttt{src/data\_sources/foresight.py} использовал \texttt{market.id} (numeric) как join key, а \texttt{src/inverse/loader.py} (который загружал исторические trades) использовал `conditionId` (hash). 99% сигналов молчал отпадал, потому что не было match'а.

\subsubsection*{Что пробовали}
\begin{enumerate}
\item Посмотреть, что возвращает Polymarket API? Оба поля возвращаются.
\item Какое поле используется в других частях кода? Нашли inconsistency.
\item Добавить логирование для match rate? Обнаружили 0% matches.
\end{enumerate}

\subsubsection*{Почему не сработало}
\textbf{Нет documentation in code comments.} Ни один из двух IDs не был задокументирован с пояснением, что один локальный, другой глобальный. Разработчик справедливо выбрал первый встреченный `id` поле.

\subsubsection*{Что сделали}
v0.9.3 (CHANGELOG, стр. 96--97):

\begin{enumerate}
\item \texttt{src/data\_sources/foresight.py}: extract `conditionId` вместо `id`.
\item \texttt{src/agents/collectors/foresight\_collector.py}: use `conditionId` как join key для matching с trade data.
\item Добавлены comments:
\begin{verbatim}
# conditionId is globally unique (CLOB + Data API)
# id is local to Gamma API (marketplace-specific)
# Always join on conditionId, not id.
\end{verbatim}
\end{enumerate}

\subsubsection*{Урок}
\textbf{Когда third-party API имеет несколько IDs, документируй purpose каждого.}

\begin{enumerate}
\item Add comment в код: ``// market.id: local (Gamma), market.conditionId: global (CLOB)''.
\item Test: fetch market через оба endpoints, verify оба ID'а return.
\item Unit test: join на wrong ID должен вернуть 0 matches (catch этот bug).
\item Добавь assertion: \texttt{assert len(matches) > 0, "No markets matched! Check join keys."}
\end{enumerate}

---

\subsection{Case Study 11: Date Serialization Crash (v0.9.5)}

\subsubsection*{Контекст}
v0.9.5 обновила timeline schemas (Phase 6). \texttt{predicted\_timeline} содержит \texttt{predicted\_date} и \texttt{target\_date} поля. Pipeline прошла all 9 stages успешно, но при попытке сохранить результат в SQLite, worker упал с:

\begin{verbatim}
TypeError: Object of type date is not JSON serializable
\end{verbatim}

Результат не был сохранён --- user смотрит на пустую страницу despite 40 минут computation (\$5--15 в cost).

\subsubsection*{Что пробовали}
\begin{enumerate}
\item Проверить, что dates парсятся корректно? Да, они \texttt{datetime.date} объекты.
\item Где serialization происходит? В worker при \texttt{model\_dump()} перед `INSERT`.
\item Почему \texttt{model\_dump()} не обрабатывает date объекты? По дефолту, \texttt{model\_dump()} не имеет \texttt{mode="json"}, который преобразует non-JSON types.
\end{enumerate}

\subsubsection*{Почему не сработало}
\textbf{Pydantic v2 breaking change.} В Pydantic v2, \texttt{model\_dump()} требует явный \texttt{mode="json"} для JSON serialization. По дефолту, это просто возвращает Python объекты. Тогда \texttt{json.dumps()} fails на `date` объектах.

\subsubsection*{Что сделали}
v0.9.5 (CHANGELOG, стр. 35--36):

\begin{enumerate}
\item Добавлены \texttt{@field\_serializer} decorators на schemas с date полями:
\begin{verbatim}
@field_serializer('predicted_date', 'target_date', 'generated_at')
def serialize_dates(self, value):
    if isinstance(value, date):
        return value.isoformat()
    return value
\end{verbatim}

\item Альтернативно: используй \texttt{model\_dump(mode="json")} везде, где нужна JSON serialization.
\item Добавлены unit tests для timeline serialization.
\end{enumerate}

\subsubsection*{Урок}
\textbf{Serialization должна быть unit-tested. Никогда не предполагай, что \texttt{model\_dump()} даст тебе JSON-ready dict.}

\begin{enumerate}
\item Add test: \texttt{model\_dump() $\to$ json.dumps() $\to$ json.loads()} round-trip.
\item При добавлении новых fields типа date/datetime, immediately add serializer.
\item Используй \texttt{model\_dump(mode="json")} в production code (более явно).
\item Помни: Pydantic v2 требует explicit `mode` для JSON types.
\end{enumerate}

---

\subsection{Case Study 12: PromptParseError Silently Dropped (v0.5.2--v0.9.4)}

\subsubsection*{Контекст}
v0.7.0 agent (EventTrendAnalyzer) отправляет JSON request на LLM. LLM (Gemini Flash) может вернуть truncated JSON (если finish\_reason=``length''). Код:

\begin{verbatim}
try:
    data = json.loads(response.text)
    return EventThreads.model_validate(data)
except (json.JSONDecodeError, ValidationError):
    return {}  # silent fallback
\end{verbatim}

Проблема: при invalid JSON, код возвращает \texttt{{}} (пустой dict). Downstream код ожидает `EventThreads` (которая имеет required fields), получает \texttt{{}} $\to$ ValidationError $\to$ assessment молчит падает, и timeline становится пустой.

\subsubsection*{Что пробовали}
\begin{enumerate}
\item Логировать ошибку? Было, но WARNING уровень, который обрезался в production.
\item Есть ли fallback logic? Да, но fallback неполный (возвращает {} вместо default EventThreads).
\item Как это вообще прошло testing? Тесты использовали MockLLMClient, который всегда возвращает valid JSON.
\end{enumerate}

\subsubsection*{Почему не сработало}
\textbf{Exception swallowing + incomplete fallback.} Catch-all `except` absorbs both JSON errors и Validation errors. Fallback \texttt{{}} не имеет структуры, которую ожидает downstream. Это как ``promise'' вернуть результат, но фактически вернуть None.

\subsubsection*{Что сделали}
v0.7.1 (CHANGELOG, стр. 275): Resilient event\_identification added.

\begin{enumerate}
\item Distinguish между JSON parse errors и validation errors:
\begin{verbatim}
try:
    data = json.loads(response.text)
    return EventThreads.model_validate(data)
except json.JSONDecodeError as e:
    logger.error("Truncated JSON: %s", response.text[:100])
    # fallback to raw headlines?
    return _fallback_event_threads(headlines)
except ValidationError as e:
    logger.error("Validation error: %s", e)
    # return empty but structured default
    return EventThreads(threads=[])
\end{verbatim}

\item Вместо silent drop, либо:
  \begin{enumerate}
  \item Re-raise exception для explicit retry
  \item Return default объект (e.g., \texttt{EventThreads(threads=[])}) со структурой
  \item Fallback to raw headlines (e.g., каждый headline = один thread)
  \end{enumerate}
\end{enumerate}

\subsubsection*{Урок}
\textbf{Never silently swallow exceptions. Either fail-fast или graceful fallback с logged reason.}

\begin{enumerate}
\item Distinguish error types: parse errors vs validation errors имеют разные fix strategies.
\item Fallback value должен иметь корректную структуру (не \texttt{{}}, а \texttt{Model(field1=[], field2=None)}).
\item Логируй exception content (e.g. \texttt{str(e)} или \texttt{e.json()}), не только error message.
\item Unit test: add case where LLM returns truncated JSON, verify fallback works.
\end{enumerate}

---

\subsection{Case Study 13: BudgetTracker Race Condition (v0.7.1)}

\subsubsection*{Контекст}
v0.7.0 реализовала budget tracking для ограничения LLM costs (\$50 default max). BudgetTracker:

\begin{verbatim}
async def check_budget(self, cost_usd: float):
    if self.total_cost + cost_usd > self.max_budget:
        raise BudgetExceededError(...)
    self.total_cost += cost_usd
\end{verbatim}

Проблема: в async context с несколькими parallel agents, два агента могут одновременно вызвать \texttt{check\_budget()}. Race condition:
\begin{enumerate}
\item Agent 1 reads \texttt{self.total\_cost = 48}.
\item Agent 2 reads \texttt{self.total\_cost = 48}.
\item Agent 1 adds 1.5, writes \texttt{self.total\_cost = 49.5}.
\item Agent 2 adds 1.5, writes \texttt{self.total\_cost = 50.0} (мог бы быть 51, что exceeds budget).
\end{enumerate}

Результат: budget check может быть bypassed.

\subsubsection*{Что пробовали}
\begin{enumerate}
\item Запустить с parallel agents? Локально, single-threaded, не воспроизводилось.
\item Добавить sleep между LLM calls? Случайно помогло в тесте, но не fix.
\item Есть ли глобальный lock? Нет.
\end{enumerate}

\subsubsection*{Почему не сработало}
\textbf{TOCTOU (Time-of-Check-Time-of-Use) race condition.} Check и Use не атомичны. В async Pythone, `await` может переключить контекст между check и write.

\subsubsection*{Что сделали}
v0.7.1 (CHANGELOG, стр. 263):

\begin{verbatim}
class BudgetTracker:
    def __init__(self):
        self._lock = asyncio.Lock()

    async def check_budget(self, cost_usd: float):
        async with self._lock:
            if self.total_cost + cost_usd > self.max_budget:
                raise BudgetExceededError(...)
            self.total_cost += cost_usd
\end{verbatim}

Теперь check и write protected by lock. Не может быть race.

\subsubsection*{Урок}
\textbf{Shared state в async code требует synchronization primitives.}

\begin{enumerate}
\item Используй \texttt{asyncio.Lock()} для защиты read-modify-write операций.
\item Никогда не предполагай, что single-threaded code safe в async context.
\item Unit test: запусти с \texttt{asyncio.gather()} несколько concurrent calls, verify atomicity.
\item Code review: флаг на любое \texttt{self.field = ...} в async function --- может быть race condition.
\end{enumerate}

---

\subsection{Case Study 14: Timeout Cascade (v0.9.4)}

\subsubsection*{Контекст}
Дефолтный timeout агентов составлял 300 секунд. При работе с тяжёлыми промптами (Opus на сложных multi-agent медиациях) LLM-ответы регулярно превышали 300s. Результат: каскадный отказ --- таймаут одного агента вызывал падение всей стадии, а с ней терялась вся предыдущая работа pipeline.

\subsubsection*{Что сделали}
Дефолтный timeout поднят с 300s до 600s (commit c96ed67). Индивидуальные таймауты настроены per stage:

\begin{enumerate}
\item \texttt{outlet\_historian}: снижен до 300s (легковесная стадия, не требует длинных LLM-ответов).
\item \texttt{delphi\_r2}: увеличен до 900s (сложная multi-agent медиация с 5 персонами, самый длинный промпт в pipeline).
\item Остальные стадии: 600s (2x от наблюдаемого p95).
\end{enumerate}

\subsubsection*{Урок}
\textbf{Дефолтные таймауты должны базироваться на 2x от наблюдаемого p95 latency, а не на оптимистичных средних.}

\begin{enumerate}
\item Мониторь реальные latency перед выставлением порогов.
\item Настраивай таймауты per stage --- одна стадия может быть 10x тяжелее другой.
\item Каскадный отказ от таймаута хуже, чем дополнительные 5 минут ожидания.
\end{enumerate}

---

\subsection{Case Study 15: Эволюция max\_tokens (v0.5.1--v0.9.5)}

\subsubsection*{Контекст}
Параметр \texttt{max\_tokens} прошёл 4 итерации:
\begin{enumerate}
\item \textbf{4096} (начальное значение) --- обрезал выходы Delphi R1, где 5 персон генерируют по 200+ токенов каждая.
\item \textbf{8192} --- всё ещё недостаточно для полных assessment'ов с обоснованиями.
\item \textbf{16384} --- обнаружилось, что OpenRouter резервирует \texttt{max\_tokens} из кредитного баланса. Установка 16384 резервировала \$5+ за вызов, даже если фактический выход --- 2000 токенов.
\item \textbf{unlimited} (текущее значение) --- cap снят полностью.
\end{enumerate}

\subsubsection*{Что сделали}
Удалён cap \texttt{max\_tokens} полностью. Ключевое открытие --- механизм резервирования кредитов OpenRouter: при \texttt{max\_tokens=16384} система блокировала эквивалент полного output на балансе пользователя, создавая иллюзию высоких расходов и блокируя параллельные вызовы при низком балансе.

\subsubsection*{Урок}
\textbf{Модели ценообразования LLM API неочевидны.} Тестируй cost impact параметра \texttt{max\_tokens} до production deploy.

\begin{enumerate}
\item Проверяй, как провайдер обрабатывает \texttt{max\_tokens}: billing по фактическому output или по зарезервированному?
\item Для OpenRouter: unlimited \texttt{max\_tokens} дешевле, чем высокий фиксированный cap.
\item Документируй pricing gotchas для каждого LLM-провайдера в проектной wiki.
\end{enumerate}

---

\subsection{Case Study 16: Инкрементальное сохранение pipeline (v0.9.5)}

\subsubsection*{Контекст}
Pipeline сохранял результаты только по полному завершению всех 9 стадий. Если QualityGate (stage 9) падал по таймауту после 40+ минут обработки, все результаты предыдущих стадий (\$5--15 в LLM costs) терялись безвозвратно. Пользователь видел пустую страницу.

\subsubsection*{Что сделали}
Реализовано per-stage инкрементальное сохранение (commit f4d0848):

\begin{enumerate}
\item Каждая завершённая стадия немедленно persists результат в \texttt{PipelineStep.output\_data}.
\item При падении на стадии $N$, стадии $1 \ldots N{-}1$ уже сохранены и доступны пользователю.
\item Worker при рестарте может продолжить с последней сохранённой стадии (checkpoint recovery).
\end{enumerate}

\subsubsection*{Урок}
\textbf{Долгоживущие pipeline обязаны делать checkpoint после каждого дорогого шага. ``Всё или ничего'' неприемлемо, когда каждый шаг стоит доллары.}

\begin{enumerate}
\item Стоимость потерянной работы пропорциональна количеству стадий без checkpoint'ов.
\item Checkpoint recovery позволяет retry с последней точки, экономя и время, и деньги.
\item Паттерн: каждый stage записывает \texttt{output\_data} сразу после успешного завершения.
\end{enumerate}

---

\section{Отложенные направления}

\subsection{Case Study 17: Domain-Specific Brier Scores}

\subsubsection*{Идея}
Разные типы рынков имеют разные точность-характеристики. Например, crypto markets более волатильны, политические markets более точны. Можно ли compute per-domain Brier Score weights и используй их для adaptive extremizing?

\subsubsection*{Почему отложено}
\textbf{Data sparsity.} из 348K informed бетторов, только ~100 имеют >5 resolved bets в each domain (crypto, politics, sports, etc.). Brier Score на малом N имеет огромную дисперсию. При попытке stratify по domains, 99% бетторов drop out из-за недостаточного разрешения данных.

Документация (docs/methodology-inverse-problem.md, стр. 398): ``Domain-specific BS --- data sparsity (1--3\% BSS gain, not worth effort)''

\subsubsection*{Урок}
\textbf{Не добавляй complexity для marginal gains на limited data.} 1--3\% улучшение требует 10x больше данных для статистической significance.

---

\subsection{Case Study 18: Bettor-Level News Correlation}

\subsubsection*{Идея}
Корреляция между news events (GDELT) и bettor positions. Гипотеза: informed бетторы быстро реагируют на breaking news, можно ли построить predictive model?

\subsubsection*{Почему отложено}
\textbf{Спекулятивная гипотеза требует специального pipeline.} Для отладки нужно:
\begin{enumerate}
\item GDELT API, которая работает (исправлена в v0.9.4, но limited к 3-месячному окну).
\item RSS pipeline для индивидуальных outlet'ов (NewsScout есть, но не per-bettor).
\item Temporal alignment: news timestamp vs trade timestamp (сложная временная модель).
\item Validation: есть ли реальная correlation? Не known.
\end{enumerate}

Документация (docs/methodology-inverse-problem.md, стр. 399): ``Bettor-level корреляция с новостями --- требует GDELT/RSS pipeline, спекулятивная гипотеза''

\subsubsection*{Урок}
\textbf{Research hypotheses требуют pilot data перед full engineering.} Не начинай building без evidence, что ideá заслуживает effort.

---

\subsection{Case Study 19: Hierarchical Trader Belief Models}

\subsubsection*{Идея}
Bettor positions отражают их beliefs. Можно ли build hierarchical Bayesian model где бетторы cluster по beliefs, и используй это для better aggregate forecast?

\subsubsection*{Почему отложено}
\textbf{Pure research project, не engineering task.} Это требует:
\begin{enumerate}
\item Новых statistical methods (Bayesian hierarchical modeling).
\item Publication-quality validation (peer review, reproducibility).
\item Большого датасета для training (мы имеем Polymarket, но это single market type).
\item Четкого определения ``belief'' в контексте prediction markets.
\end{enumerate}

Документация (docs/methodology-inverse-problem.md, стр. 400): ``Иерархические модели --- research project, не engineering task''

\subsubsection*{Урок}
\textbf{Distinguish между engineering (solve known problem) и research (formulate new problem).} Не заполняй product roadmap research ideas.

---

\subsection{Case Study 20: Kalshi API Integration}

\subsubsection*{Идея}
Kalshi --- второй major US prediction market (после Polymarket). Может ли Delphi интегрировать Kalshi для расширения охвата?

\subsubsection*{Почему отложено}
\textbf{US-only coverage, low ROI.} Kalshi имеет:
\begin{enumerate}
\item ~500 active markets (vs Polymarket ~5000+).
\item Regional exposure (США только).
\item Higher regulation (SEC oversight).
\item No public bettor profiles (в отличие от Polymarket, который имеет on-chain transparency).
\end{enumerate}

Cost-benefit: интеграция требует ~3 дня (API integration, schema changes, tests). Результат: +200 рынков (low quality), 0 additional signals (no public bettors). Not worth.

\subsubsection*{Урок}
\textbf{Evaluate ROI перед adding new data source.} Не все APIs стоят интегрирования.

---

\subsection{Case Study 21: BigQuery for GDELT Historical Data}

\subsubsection*{Идея}
GDELT имеет полный historical dataset в BigQuery (2.65 TB/year). Можно ли использовать для batch retrospective testing?

\subsubsection*{Почему отложено}
\textbf{Cost-prohibitive at scale.} BigQuery pricing:
\begin{enumerate}
\item 1 TB free/месяц, затем \$6.25/TB.
\item Unoptimized query: 311 GB scanned = \$1.94 per query.
\item Optimized query с temporal decorators: 6.28 GB scanned = \$0.04 per query.
\item Для 100-query batch retrospective: \$400--\$200 (depends on optimization).
\end{enumerate}

Alternative: 15-minute CSV polling (free, 2--5 MB per batch) + local DuckDB. Total cost: \$0. Speed: slower, но достаточна для production monitoring.

Документация (\texttt{tasks/research/gdelt\_api.md}, стр. 43--47): ``BigQuery viable для retrospective, но для real-time --- используй CSV polling''

\subsubsection*{Урок}
\textbf{Evaluate cloud costs vs local alternatives. Free polling beats paid queries.}

---

\section*{Summary Table: Dead Ends and Lessons}

\begin{table}[h]
\centering
\small
\begin{tabular}{|l|l|l|l|l|l|}
\hline
\textbf{\#} & \textbf{Category} & \textbf{Case Study} & \textbf{Status} & \textbf{Version} & \textbf{Impact} \\
\hline
1 & \textbf{API} & Metaculus 403 & Fixed, Disabled & v0.5.1, v0.9.4 & Auth-required tier lock \\
2 &  & GDELT Cyrillic & Fixed & pre-v0.5.1, v0.9.4 & HTML crashes prevented \\
3 &  & Polymarket camelCase & Fixed & pre-v0.5.1 & 422 error resolved \\
4 &  & Reuters RSS & Removed & v0.9.4 & 404 dead feed \\
\hline
5 & \textbf{Architecture} & YandexGPT & Removed & v0.8.0 & Consolidated to OpenRouter \\
6 &  & Sonnet 4.6 & Removed & v0.9.4 & Non-existent model \\
7 &  & Dark mode & Removed & v0.8.0 & Unmaintained feature \\
8 &  & Pico.css $\to$ Tailwind & Migrated & v0.8.0 & Classless CSS limit \\
\hline
9 & \textbf{Critical Bugs} & Temporal leak & Fixed & v0.9.2 & BSS +0.092$\to$+0.127 \\
10 &  & conditionId & Fixed & v0.9.3 & 99\% signal loss \\
11 &  & Date serialization & Fixed & v0.9.5 & Complete data loss \\
12 &  & PromptParseError & Fixed & v0.5.2--v0.9.4 & Silent assessment drop \\
13 &  & BudgetTracker race & Fixed & v0.7.1 & Budget bypass \\
14 &  & Timeout cascade & Fixed & v0.9.4 & Cascading stage failures \\
15 &  & max\_tokens evolution & Fixed & v0.5.1--v0.9.5 & Credit reservation bloat \\
16 &  & Incremental save & Fixed & v0.9.5 & \$5--15 result loss \\
\hline
17 & \textbf{Deferred} & Domain-specific BS & Deferred & --- & 1--3\% gain, high sparsity \\
18 &  & Bettor-news correlation & Deferred & --- & Speculative hypothesis \\
19 &  & Hierarchical beliefs & Deferred & --- & Pure research \\
20 &  & Kalshi API & Deferred & --- & US-only, low ROI \\
21 &  & BigQuery GDELT & Deferred & --- & Cost-prohibitive (\$0.04--\$1.94/query) \\
\hline
\end{tabular}
\end{table}

\noindent
\textbf{Key Takeaways:}
\begin{enumerate}
\item \textbf{External APIs require continuous monitoring.} Breaking changes happen (Metaculus, Reuters). Subscribe to official channels.
\item \textbf{Single-provider architecture better than unreliable fallbacks.} YandexGPT was removed, OpenRouter is sole LLM provider.
\item \textbf{Serialization and temporal correctness are non-negotiable.} Date serialization crash and temporal leak both caused complete data loss / model degradation.
\item \textbf{Distinguish engineering from research.} Hierarchical models are research; bettor signals engineering.
\item \textbf{ROI analysis prevents feature creep.} Domain-specific BS, Kalshi, BigQuery: all deferred because cost > benefit.
\end{enumerate}


% =====================================================================
% PART X: Roadmap
% =====================================================================

\part*{X. Дорожная карта}

\section{Ближайшие задачи}

\begin{enumerate}
  \item \textbf{Market Dashboard --- live data pipeline.}
    Shell-cron refresh профилей (v0.9.3) уже работает.
    Осталось: ARQ-based incremental rebuild + hot-reload в \texttt{MarketSignalService}.
    Критерий: $\geq 10$ рынков с informed consensus.

  \item \textbf{Pipeline Checkpoint \& Resume.}
    Сериализация \texttt{PipelineContext} после каждой стадии
    в \texttt{PipelineStep.output\_data} (поле существует, но не заполняется).
    Resume: \texttt{POST /api/v1/predictions/\{id\}/resume}.
    Оценка: ${\sim}200{-}300$ LOC, без миграции БД.

  \item \textbf{Gemini Flash --- broken JSON.}
    Light-пресет генерирует кривые JSON на ряде стадий.
    Варианты: JSON repair middleware, prompt engineering,
    замена на Claude Haiku 4.5.

  \item \textbf{Event-level prediction storage.}
    Сохранять \texttt{PredictedTimeline} в JSON/DB после каждого run.
    Блокирует per-prediction market eval (сравнение Delphi BS vs.\ Polymarket BS для конкретных прогнозов).
    Агрегатный market eval (направления B+C) уже реализован.

  \item \textbf{Пилот ретроспективной оценки.}
    Инфраструктура готова (Brier, Log Score, Composite Score, bootstrap CI, Wayback CDX).
    Осталось: 50 runs $\times$ 3 горизонта $\to$ ${\sim}150{-}350$ пар.
    Стоимость $< \$1$.
\end{enumerate}

\section{Backlog}

\begin{table}[h]
\centering
\begin{tabular}{clll}
\toprule
\textbf{\#} & \textbf{Задача} & \textbf{Приоритет} & \textbf{Зависимости} \\
\midrule
B.2  & Kalshi, OECD integration         & Низкий   & --- \\
B.4  & Калибровка BERTScore порогов      & Средний  & После пилота eval \\
B.5  & Platt scaling в Judge             & Средний  & Reliability $> 0.05$ \\
B.6  & Per-persona weight update (Brier) & Средний  & Per-persona BS разброс $> 0.10$ \\
B.9  & Правовая проверка licensing       & Средний  & До публичного деплоя \\
B.10 & Event-level storage (JSON/DB)     & Высокий  & Блокирует market eval \\
B.11 & HuggingFace dataset для N=500+    & Низкий   & --- \\
\bottomrule
\end{tabular}
\caption{Backlog задач после production deploy\protect\footnotemark}
\label{tab:backlog}
\end{table}
\footnotetext{Полный backlog содержит 10 задач; показаны 7 приоритетных.}

\section{Дискуссия: архетипы стратегий принятия решений}
\label{sec:discussion_archetypes}

Раздел фиксирует тезисы рабочей сессии команды (04.04.2026),
посвящённой расширению модуля обратной задачи (Часть~VI)
за счёт стратегической классификации участников рынка.

\subsection{Центральная гипотеза}

Участники рынков предсказаний принимают решения
в рамках конечного множества \emph{архетипов стратегий},
описанных в литературе по теории принятия решений,
поведенческой экономике и бизнес-аналитике.
Если по истории торговли отдельного актора можно определить его архетип,
то появляется возможность \emph{более точно оценивать bias}
каждой группы участников --- переоценку или недооценку событий ---
и использовать эту информацию для коррекции informed consensus.

\subsection{Ключевые тезисы}

\begin{enumerate}
  \item \textbf{Правила внутри правил.}
    Сам механизм принятия решений меняется во времени:
    не только входные данные, но и правила их обработки.
    На малых горизонтах актор следует фиксированной стратегии (policy),
    на больших --- возможна смена стратегии.
    Мета-правила этих смен --- отдельная исследовательская задача.

  \item \textbf{Извлечение архетипов из литературы.}
    LLM обрабатывает корпус литературы по принятию решений
    (бизнес-методологии, когнитивная психология, трейдинг,
    «Искусство войны», методички хедж-фондов)
    и извлекает ${\sim}20{-}30$ архетипов:
    высокочастотный скальпер, фундаментальный аналитик,
    реактор на новости, контрариан и т.\,д.

  \item \textbf{LLM как soft-классификатор.}
    Для каждого актора с известной историей торговли
    (timestamp, buy/sell, volume, price) + внешние данные (новости, отчёты)
    LLM оценивает: «если актор X следует архетипу A,
    согласуется ли это с его реальными действиями?»
    Байесовское правдоподобие по всем архетипам
    даёт мягкую классификацию.

  \item \textbf{Обучаемые графы решений.}
    Каждый архетип формализуется как граф:
    вершины = входные данные и промежуточные решения,
    рёбра = правила вывода с весами.
    Локальные перестройки графа (добавление вершин, изменение весов,
    отсечение логических противоречий) --- аналог
    дискретного градиентного спуска --- оптимизируют
    одновременно качество кластеризации и точность прогноза.

  \item \textbf{Скрытые переменные.}
    История торговли --- лишь часть переменных;
    мотивация, нагугленная информация, личный опыт --- скрыты.
    Архетипы --- попытка восстановить скрытые переменные
    через структурные предположения из литературы
    (в отличие от чисто data-driven кластеризации).

  \item \textbf{Bias по архетипам.}
    Зная архетип трейдеров на конкретном рынке,
    можно определить: \emph{переоценивают} они событие или \emph{недооценивают}.
    Пример: если на рынке олова доминируют «реакторы на новости»,
    а аналитический отчёт ещё не вышел --- рынок может недооценивать событие.

  \item \textbf{Суперпрогнозисты vs случайность.}
    Прежде чем доверять «суперпрогнозистам»,
    необходимо сравнить их результат с нулевой гипотезой:
    при $N$ участниках ${\sim}k$ из них будут в плюсе
    просто по закону больших чисел.
    Bayesian shrinkage (Часть~VI, $k=15$) частично решает эту проблему,
    но явный статистический тест (permutation test или bootstrap)
    был бы строже.

  \item \textbf{Per-persona / per-market анализ.}
    Может оказаться, что разные Дельфи-персоны
    (Реалист, Геостратег, Экономист, Медиа-эксперт, Адвокат дьявола)
    систематически лучше прогнозируют на разных типах рынков.
    Это направление пересекается с backlog-задачей B.6
    (per-persona weight update по Brier Score).
\end{enumerate}

\subsection{Соотношение с текущей системой}

\begin{table}[ht]
\centering
\small
\begin{tabular}{p{5cm}p{5cm}p{4.5cm}}
\toprule
\textbf{Тезис} & \textbf{Что есть в Delphi Press} & \textbf{Разрыв} \\
\midrule
Кластеризация трейдеров по архетипам &
  HDBSCAN по Brier Score + volume + timing &
  Кластеризация по \emph{точности}, не по \emph{стратегии} \\
LLM извлекает архетипы &
  5 фиксированных персон (промпты) &
  Персоны для генерации прогнозов, не для классификации трейдеров \\
Bias estimation по архетипу &
  Informed consensus (BSS $+0.196$) &
  Агрегация по точности, не по типу мышления \\
Обучаемые графы &
  Нет &
  Новое направление R\&D \\
Суперпрогнозисты vs случайность &
  Bayesian shrinkage ($k=15$) &
  Нет явного permutation test \\
Per-persona weight &
  Статичные initial\_weight &
  Backlog B.6 \\
\bottomrule
\end{tabular}
\caption{Тезисы дискуссии vs.\ текущая реализация}
\label{tab:discussion_gap}
\end{table}

\subsection{Оценка трудоёмкости и приоритеты}

\begin{table}[ht]
\centering
\small
\begin{tabular}{p{5.5cm}ccp{4cm}}
\toprule
\textbf{Направление} & \textbf{Сложность} & \textbf{Оценка} & \textbf{Зависимости} \\
\midrule
Извлечение архетипов из литературы (LLM research) &
  Средняя & 1--2 нед. & Корпус литературы \\
Расширение bettor features (frequency, reaction latency) &
  Низкая & 3--5 дней & Расширение \texttt{profiler.py} \\
LLM-классификация top-100 INFORMED &
  Средняя & 1--2 нед. & Архетипы + features \\
Bias estimation по архетипам &
  Высокая & 2--3 нед. & Качество классификации \\
Обучаемые графы стратегий &
  Очень высокая & 2--3 мес. & Новая архитектура \\
Мета-правила (правила смены правил) &
  Исследоват. & 6+ мес. & Фундаментальный R\&D \\
\bottomrule
\end{tabular}
\caption{Оценка трудоёмкости направлений из дискуссии}
\label{tab:discussion_effort}
\end{table}

\subsection{Рекомендуемые следующие шаги}

\begin{enumerate}
  \item \textbf{Proof of concept}: извлечь архетипы из 10--15 источников,
    провести LLM-классификацию на top-100 INFORMED кошельках,
    проверить корреляцию архетипа с Brier Score.
  \item \textbf{Расширение profiler}: добавить trade frequency,
    avg volume, reaction latency к \texttt{BettorProfile}.
  \item \textbf{Permutation test}: для INFORMED tier ---
    явный тест против нулевой гипотезы (случайная выборка
    из $N$ участников даёт $k$ ``суперпрогнозистов'').
  \item \textbf{Per-persona analysis}: запустить бэктест
    с отдельным tracking Brier Score каждой Дельфи-персоны
    (блокируется задачей B.6).
\end{enumerate}

% =====================================================================
% APPENDICES
% =====================================================================

\begin{appendices}

% -- Appendix A: Data Flow ---------------------------------------------
\section{Data Flow: 9 стадий pipeline}
\label{app:dataflow}

\begin{table}[h]
\centering
\footnotesize
\begin{tabular}{clp{3.5cm}lccl}
\toprule
\textbf{\#} & \textbf{Стадия} & \textbf{Агенты} & \textbf{Par?} & \textbf{min} & \textbf{Timeout} & \textbf{Output} \\
\midrule
1 & Collection      & NewsScout, EventCalendar, OutletHistorian, Foresight & Да & 2/4 & 600s & signals, events, profile \\
2 & Event ID        & EventTrendAnalyzer      & Нет & ---  & 600s & 20 EventThread \\
3 & Trajectory      & Geo, Econ, Media        & Да  & 2/3  & 600s & assessments \\
4 & Delphi R1       & 5 персон                & Да  & 3/5  & 600s & R1 assessments \\
5 & Delphi R2       & Mediator $\to$ 5 персон & Mix & 3/5  & 900s & synthesis + R2 \\
6 & Consensus       & Judge                   & Нет & ---  & 300s & timeline + top-7 \\
7 & Framing         & FramingAnalyzer         & Нет & ---  & 300s & 7 briefs \\
8 & Generation      & StyleReplicator         & Нет & ---  & 300s & 14-21 headlines \\
9 & Quality Gate    & QualityGate             & Нет & ---  & 300s & 5-7 final \\
\bottomrule
\end{tabular}
\caption{Pipeline: стадии, агенты, параметры}
\label{tab:pipeline_stages}
\end{table}

% -- Appendix B: Prompts (from subagent 7) ----------------------------
\section{Промпты агентов}
\label{app:prompts}

Ниже приведены полные тексты системных промптов всех агентов
Дельфи-пайплайна. Промпты хранятся в директории
\texttt{docs/prompts/} и передаются LLM-моделям через OpenRouter API.

\subsection{Реалист (Realist)}
\label{app:prompt_realist}
Файл: \texttt{docs/prompts/realist.md} (397 строк).
Аналитическая рамка: базовые ставки, исторические прецеденты,
институциональная инерция (Tetlock 2005). Начальный вес: 0.22.

\subsection{Геополитический стратег (Geostrateg)}
\label{app:prompt_geostrateg}
Файл: \texttt{docs/prompts/geostrateg.md} (459 строк).
Аналитическая рамка: неореализм Уолца + конструктивизм,
cui bono, деревья решений. Начальный вес: 0.20.

\subsection{Экономический аналитик (Economist)}
\label{app:prompt_economist}
Файл: \texttt{docs/prompts/economist.md} (471 строка).
Аналитическая рамка: follow the money, рациональный актор с бюджетными ограничениями,
экономический календарь. Начальный вес: 0.20.

\subsection{Медиа-эксперт (Media Expert)}
\label{app:prompt_media_expert}
Файл: \texttt{docs/prompts/media-expert.md} (471 строка).
Аналитическая рамка: гейткипинг (White 1950), фрейминг (Entman 1993),
6 критериев новостной ценности. Начальный вес: 0.18.

\subsection{Адвокат дьявола (Devil's Advocate)}
\label{app:prompt_devils_advocate}
Файл: \texttt{docs/prompts/devils-advocate.md} (503 строки).
Аналитические инструменты: pre-mortem, steelmanning, чёрные лебеди (Талеб).
Цель --- не минимизация BS, а генерация контраргументов. Начальный вес: 0.20.

\subsection{Медиатор (Mediator)}
\label{app:prompt_mediator}
Файл: \texttt{docs/prompts/mediator.md} (435 строк).
Задача: классифицировать результаты R1 (consensus/disputes/gaps),
сформулировать ключевые вопросы, анонимизировать позиции (Expert A-E).

\subsection{Судья-арбитр (Judge)}
\label{app:prompt_judge}
Файл: \texttt{docs/prompts/judge.md} (639 строк).
Алгоритм: weighted median + Platt scaling + headline selection + wild cards.
Horizon-adaptive weights. Детерминистический (без LLM с v0.7.0).

\bigskip
\noindent\emph{Полные тексты промптов доступны в репозитории по указанным путям.
Они слишком объёмны для включения в печатный документ (суммарно ${\sim}3400$ строк).}


% -- Appendix C: Known Gotchas ----------------------------------------
\section{Известные архитектурные ограничения}
\label{app:gotchas}

\begin{enumerate}[nosep]
  \item \texttt{Dict vs Pydantic}: слоты контекста типизированы как \texttt{Any};
    потребители должны обрабатывать оба формата.
  \item \texttt{max\_tokens=4096}: недостаточно для R1 Delphi (5-15 items $\times$ 200 tokens);
    может обрезать (\texttt{finish\_reason="length"}).
  \item \texttt{EventTrendAnalyzer} заполняет 3 слота:
    логически Stage 2, но выдаёт trajectories (Stage 3).
  \item Stage 5 \texttt{StageDefinition} списывает только mediator;
    \texttt{\_run\_delphi\_r2()} содержит custom logic для 5 персон.
  \item Round detection через \texttt{mediator\_synthesis}:
    если медиатор упал, R2 персоны думают, что это R1.
  \item \texttt{ForesightCollector} не вызывает LLM:
    pure data agent, но conformant к \texttt{BaseAgent}.
  \item \texttt{StyleReplicator}: выбор модели по языку издания;
    если язык определён неверно --- неправильный task ID.
  \item Single LLM provider (OpenRouter):
    нет fallback на альтернативного провайдера.
  \item Budget \$50 default:
    один Opus call стоит \$5-15; два прогноза могут исчерпать бюджет.
  \item Persona weights статичны:
    Judge использует \texttt{initial\_weight} с horizon-поправками,
    но Brier-обновление весов не реализовано (backlog B.6).
  \item QualityGate REVISE$\to$drop:
    revision не реализована; REVISE = немедленный drop (без попыток).
\end{enumerate}

\end{appendices}

% =====================================================================
% Bibliography (will be merged from all parts)
% =====================================================================

\begin{thebibliography}{99}

\bibitem[Manski(2006)]{manski2006interpreting}
Manski, C.~F. (2006).
\newblock Interpreting the predictions of prediction markets.
\newblock \emph{Economics Letters}, 91(3):425--429.

\bibitem[Satop\"a\"a et~al.(2014)]{satopaa2014combining}
Satop\"a\"a, V.~A., Baron, J., Foster, D.~P., Mellers, B.~A., Tetlock, P.~E., \& Ungar, L.~H. (2014).
\newblock Combining multiple probability predictions using a simple logit model.
\newblock \emph{International Journal of Forecasting}, 30(2):344--356.

\bibitem[Ferro \& Fricker(2012)]{ferro2012sampling}
Ferro, C.~A.~T. \& Fricker, T.~E. (2012).
\newblock A bias-corrected decomposition of the Brier score.
\newblock \emph{Quarterly Journal of the Royal Meteorological Society}, 138(668):1954--1960.

\bibitem[Burnham \& Anderson(2002)]{burnham2002model}
Burnham, K.~P. \& Anderson, D.~R. (2002).
\newblock \emph{Model Selection and Multimodel Inference}.
\newblock Springer, 2nd edition.

\bibitem[Ross et~al.(2011)]{ross2011reduction}
Ross, S., Gordon, G.~J., \& Bagnell, J.~A. (2011).
\newblock A reduction of imitation learning and structured prediction to no-regret online learning.
\newblock In \emph{AISTATS}, pages 627--635.

\bibitem[Sun et~al.(2023)]{sun2023mega}
Sun, X., Yang, S., \& Mangharam, R. (2023).
\newblock {MEGA-DAgger}: Imitation learning with multiple imperfect experts using topology-based performance evaluation.
\newblock \emph{arXiv preprint arXiv:2303.00638}.

\bibitem[Xu et~al.(2022)]{xu2022dwbc}
Xu, H., Zhan, X., Yin, H., \& Qin, H. (2022).
\newblock Discriminator-weighted offline imitation learning from suboptimal demonstrations.
\newblock In \emph{ICML}, pages 24725--24742.

\bibitem[Karpoora~Sundara~Pandian \& Baheri(2025)]{densityratio2025}
Karpoora~Sundara~Pandian, S. \& Baheri, A. (2025).
\newblock Density-ratio weighted behavioral cloning: Learning control policies from corrupted datasets.
\newblock \emph{arXiv preprint arXiv:2510.01479}.

\bibitem[Brown et~al.(2019)]{brown2019irl}
Brown, D.~S., Goo, W., Nagarajan, P., \& Niekum, S. (2019).
\newblock Extrapolating beyond suboptimal demonstrations via inverse reinforcement learning from observations.
\newblock In \emph{ICML}, pages 783--792.

\bibitem[Beliaev \& Pedarsani(2025)]{irleed2024}
Beliaev, M. \& Pedarsani, R. (2025).
\newblock {IRLEED}: Inverse reinforcement learning by estimating expertise of demonstrators.
\newblock In \emph{AAAI-25}, Vol.~39, pages 15532--15540.

\bibitem[Shen et~al.(2026)]{dibm2025}
Shen, W., Ma, J., Zhou, M., \& Meng, Z. (2026).
\newblock Learning diverse skills for behavior models with mixture of experts.
\newblock \emph{arXiv preprint arXiv:2601.12397}.

\bibitem[{Attention-Based BC}(2024)]{attention_bc2024}
Zhang, Y. et~al. (2024).
\newblock An attention-based behavioral cloning approach for algorithmic trading.
\newblock \emph{Applied Intelligence}, 54:15303--15324.

\bibitem[Tiapkin et~al.(2024)]{tiapkin2024demo}
Tiapkin, D., Belomestny, D., Calandriello, D., Moulines, E., Naumov, A., Perrault, P., Valko, M., \& M\'enard, P. (2024).
\newblock Demonstration-regularized {RL}.
\newblock In \emph{ICLR 2024}.

\bibitem[Akey et~al.(2025)]{akey2025polymarket}
Akey, P., Gr\'egoire, V., Harvie, C., \& Martineau, C. (2025).
\newblock Who wins and who loses in prediction markets? Evidence from {Polymarket}.
\newblock \emph{SSRN 6443103}.

\bibitem[Mitts \& Ofir(2026)]{mitts2026informed}
Mitts, J. \& Ofir, M. (2026).
\newblock From Iran to Taylor Swift: Informed trading in prediction markets.
\newblock \emph{Harvard Law Corporate Governance Forum}.

\bibitem[Clinton \& Huang(2024)]{clinton2024polymarket}
Clinton, J. \& Huang, T. (2024).
\newblock Polymarket accuracy study.
\newblock Vanderbilt University.

\bibitem[B\"urgi et~al.(2025)]{burgi2025kalshi}
B\"urgi, C., Deng, L., \& Whelan, K. (2025).
\newblock Makers and takers: The economics of the Kalshi prediction market.
\newblock \emph{CEPR Discussion Paper}.

\bibitem[Wolfers \& Zitzewitz(2004)]{wolfers2004prediction}
Wolfers, J. \& Zitzewitz, E. (2004).
\newblock Prediction markets.
\newblock \emph{Journal of Economic Perspectives}, 18(2):107--126.

\bibitem[Tetlock \& Gardner(2015)]{tetlock2015superforecasting}
Tetlock, P.~E. \& Gardner, D. (2015).
\newblock \emph{Superforecasting: The Art and Science of Prediction}.
\newblock Crown Publishers.

\end{thebibliography}

\end{document}
